\chapter{Hall fields on curves}
\label{v4-arith:curve:chapter}

A torsion sheaf on a curve splits into its parts at closed points.
Hall multiplication and the Green form split with it. This elementary
decomposition produces operators whose exchange coefficients count
points over every finite extension of the ground field. The coefficient
depends on the normalization of both operators. Keeping the Hall
creators and their Green adjoints gives a mass series. Changing their
normalization gives the curve zeta function.

We follow these operators into three recipients. Polynomial Hall
states carry local Heisenberg fields, and finite power covers retain
their residue maps. Collective counting modes instead lead to fields
with finitely many multiplicative collision divisors after reciprocal
normalization. The unchanged Green exchange has a meromorphic
difference equation and a positive coherent state. Derived intersections
of the collision divisors explain why their scalar equations and
supported differential modules require different comparisons.

Prime-indexed Macdonald spaces separate the choice of local parameters
from a geometric Hall realization. A final nilpotent-coefficient
example constructs the scalar kernel and its finite-jet comparison,
then locates the missing critical Hall map.

\section{Closed points and Green adjunction}
\label{v4-arith:curve:sec:green}

Let $C$ be a smooth projective geometrically connected curve over
$\mathbb F_q$. For a closed point $x$, write
$d_x=[k(x):\mathbb F_q]$ and $Q_x=q^{d_x}$.
A choice of parameter and coefficient field identifies the completed
local ring with $\mathbb F_{Q_x}[[t]]$. Its finite-length modules
are the finite nilpotent $\mathbb F_{Q_x}[t]$-modules.
The resulting Hall algebra is independent of this presentation.
The presentation names its Jordan partitions.

Use complex coefficients. For a prime power $Q$, let $H_Q$ have basis
$u_M$ indexed by nilpotent modules and put
$a_M=|\operatorname{Aut}M|$. The product and positive Hermitian form are
\[
 u_Mu_N=\sum_L F^L_{M,N}u_L,
 \qquad \langle u_M,u_N\rangle_G=\frac{\delta_{MN}}{a_M},
\]
where $F^L_{M,N}$ counts submodules of type $N$ with quotient $M$.
The form is conjugate-linear in its first argument.
The torsion Euler form is zero. Thus the Green coproduct is adjoint
to this product without an Euler twist.

The local primitive theorem, with its normalization fixed in
Theorem~\ref{v4-arith:hall:thm:primitives}, gives
\begin{equation}\label{v4-arith:curve:eq:primitives}
 \begin{gathered}
 H_Q=\mathbb C[p_1,p_2,\ldots],\qquad
 \Delta p_r=p_r\otimes1+1\otimes p_r,\\
 p_r=\sum_{|\lambda|=r}
       \prod_{j=1}^{\ell(\lambda)-1}(1-Q^j)u_\lambda,
 \qquad
 \langle p_r,p_s\rangle_G=\delta_{rs}\frac r{Q^r-1}.
 \end{gathered}
\end{equation}
In particular the coefficient of the indecomposable module of
length $r$ in $p_r$ is one.

The elementary generating series gives another way to see the
normalization. Put
\[
 e_r=Q^{r(r-1)/2}u_{1^r},\qquad
 E(t)=\sum_{r\geq0}e_rt^r,
 \qquad e_0=1.
\]
Subspace counting gives
$\Delta u_{1^r}=\sum_{i+j=r}Q^{-ij}u_{1^i}\otimes u_{1^j}$,
so $E$ is group-like.

Products of semisimple classes are triangular in the Jordan basis.
The successive kernel dimensions specify a filtration lowered by $t$.
Equality with these dimensions forces the kernel filtration itself. Its coefficient is
one. Other terms have larger kernel dimensions. This proves the
polynomial presentation in each finite Hall degree.

The logarithm
then gives the primitive generators through
\[
 \log E(t)=\sum_{r\geq1}\frac{(-1)^{r-1}}r p_rt^r.
\]

Since
$|\operatorname{GL}_r(\mathbb F_Q)|
=Q^{r(r-1)/2}\prod_{j=1}^r(Q^j-1)$, their norms can also be read from
\[
 \langle E(u),E(v)\rangle_G
 =\sum_{r\geq0}\frac{Q^{r(r-1)/2}(\bar uv)^r}
                         {\prod_{j=1}^r(Q^j-1)}
 =\prod_{j\geq1}(1-Q^{-j}\bar uv)^{-1}.
\]
The last equality follows from $(1-t)P(Qt)=P(t)$ and $P(0)=1$.

Primitivity and Hopf adjunction give factorial norms for powers
of a primitive. Taking logarithms gives
\eqref{v4-arith:curve:eq:primitives}.

\begin{proposition}\label{v4-arith:curve:prop:adjoints}
On the polynomial domain, multiplication $M_r$ by $p_r$ has Green
adjoint
\[
 D_r=\frac r{Q^r-1}\partial_{p_r},\qquad
 [D_r,M_s]=\delta_{rs}\frac r{Q^r-1}I.
\]
The inverse pairing, completed by nonnegative Hall weight, is
\begin{equation}\label{v4-arith:curve:eq:inverse-pairing}
 \sum_\lambda u_\lambda\otimes a_\lambda u_\lambda
 =\exp\left(\sum_{r\geq1}\frac{Q^r-1}{r}p_r\otimes p_r\right).
\end{equation}
\end{proposition}
\begin{proof}
The adjoint of multiplication by a primitive is a derivation.
Its value on $p_s$ is the scalar
$\delta_{rs}r/(Q^r-1)$, which determines the derivation on $H_Q$.
Iterating adjunction makes the monomials orthogonal, with
\[
 \left\|\prod_r p_r^{m_r}\right\|_G^2
 =\prod_r m_r!\left(\frac r{Q^r-1}\right)^{m_r}.
\]
Sum each monomial tensored with its dual to obtain the exponential.
The Hall basis gives the left side of
\eqref{v4-arith:curve:eq:inverse-pairing}. Each weight space is finite
dimensional, so both expressions define the same completed tensor.
\end{proof}

The torsion Hall space of $C$ is the algebraic tensor product
\[
 \mathcal F_C=\bigotimes_{x\in|C|}' H_{q^{d_x}},
 \qquad \deg p_{x,r}=rd_x.
\]
The prime means that a tensor is the vacuum $1$ at all but finitely
many points. Hom and extension groups between different supports
vanish, so products and Green forms factor.

There are finitely many
closed points of bounded degree. Hence each weight space is finite
dimensional.

Define, for $n>0$,
\begin{equation}\label{v4-arith:curve:eq:global-modes}
 \begin{split}
 N_n&=\sum_{d_x\mid n}d_x=\#C(\mathbb F_{q^n}),\\
 A_{-n}&=\sum_{d_x\mid n}d_xM_{x,n/d_x},\qquad
 B_n=\sum_{d_x\mid n}d_xD_{x,n/d_x}.
 \end{split}
\end{equation}
All sums are finite. A point of degree $d$ has exactly $d$ geometric
points over $\mathbb F_{q^n}$ when $d\mid n$, and none otherwise.
Operators at different points commute. Therefore
\begin{equation}\label{v4-arith:curve:eq:green-commutator}
 [B_n,A_{-m}]=\delta_{nm}\frac{nN_n}{q^n-1}I,
 \qquad B_n=A_{-n}^{*G}.
\end{equation}
Two positive modes commute, as do two negative modes.

\section{Three normalizations of the exchange}
\label{v4-arith:curve:sec:exchange}

For the unchanged creators and Green adjoints set
\[
 \Gamma_-(w)=\exp\left(\sum_{n>0}\frac{A_{-n}w^n}{n}\right),
 \qquad
 \Gamma_+^G(z)=\exp\left(\sum_{n>0}\frac{B_nz^{-n}}n\right).
\]
These are coefficientwise series in the independent formal variables
$w,z^{-1}$. Each coefficient is an operator on $\mathcal F_C$.
Creation may equivalently be evaluated in the completion by weight.
Annihilation is locally finite on every polynomial state.

\begin{theorem}\label{v4-arith:curve:thm:exchange}
Put
\[
 K_G(t)=\exp\left(\sum_{n>0}\frac{N_nt^n}{n(q^n-1)}\right),
 \qquad Z_C(t)=\exp\left(\sum_{n>0}\frac{N_nt^n}{n}\right).
\]
Then
\begin{align}
 \Gamma_+^G(z)\Gamma_-(w)
   &=K_G(w/z)\Gamma_-(w)\Gamma_+^G(z),
       \label{v4-arith:curve:eq:green-exchange}\\
 K_G(qt)&=Z_C(t)K_G(t),
       \label{v4-arith:curve:eq:difference}\\
 Z_C(t)&=\prod_x(1-t^{d_x})^{-1},\qquad
 K_G(t)=\prod_{j\geq1}Z_C(q^{-j}t).
       \label{v4-arith:curve:eq:products}
\end{align}
The last product uses convergent logarithmic coefficients at the
fixed numerical value $q>1$.
\end{theorem}
\begin{proof}
The commutator of the two exponents is the central series
$\sum_{n>0}N_n(w/z)^n/[n(q^n-1)]$.
Formal conjugation of the exponentials proves
\eqref{v4-arith:curve:eq:green-exchange}.
Subtracting the logarithms at $qt$ and $t$ proves the difference
equation. Writing $n=rd_x$ in $\log Z_C$ gives
$\sum_x\sum_{r>0}t^{rd_x}/r$, the logarithm of the Euler product.
Finally $1/(q^n-1)=\sum_{j\geq1}q^{-jn}$ proves the product for
$K_G$ in each logarithmic coefficient.
\end{proof}

There are two ways to obtain $Z_C$ itself as an exchange factor.
First keep the creators and replace $B_n$ by
$J_n=(q^n-1)B_n$. The corresponding annihilation vertex is
\begin{equation}\label{v4-arith:curve:eq:zeta-annihilator}
 \Gamma_+^\zeta(z)
 =\Gamma_+^G(z/q)\Gamma_+^G(z)^{-1}
 =\exp\left(\sum_{n>0}\frac{J_nz^{-n}}n\right).
\end{equation}
It has exchange factor $Z_C(w/z)$ with $\Gamma_-(w)$.
Its modes are no longer the Green adjoints of the fixed creators.

Second use the positive real square root and put
\[
 \widehat A_{-n}=\sqrt{q^n-1}\,A_{-n},\qquad
 \widehat B_n=\sqrt{q^n-1}\,B_n.
\]
They remain Green adjoints and satisfy
$[\widehat B_n,\widehat A_{-m}]=\delta_{nm}nN_nI$.
Their half-vertices again exchange by $Z_C(w/z)$, with both
operators changed. Locally this is the adjoint pair
\[
 a_{-r}=\sqrt{Q^r-1}\,M_r,\qquad
 a_r=\frac r{\sqrt{Q^r-1}}\partial_{p_r}.
\]

All three constructions have a common polynomial domain.
No everywhere-defined operator on the Green Hilbert completion has
been asserted. The dilation $z\mapsto z/q$ in
\eqref{v4-arith:curve:eq:zeta-annihilator} is a spectral change of
variable. Absolute Frobenius is the different map $z\mapsto z^q$.

An independent oscillator polynomial space
$\mathbb C[u_{x,r}:x\in|C|,r>0]$ is identified with $\mathcal F_C$
by $u_{x,r}\mapsto p_{x,r}$. This preserves weight, multiplication
and vacuum. It sends multiplication by $u_{x,r}$ to $M_{x,r}$
and $r\partial_{u_{x,r}}$ to $(q^{rd_x}-1)D_{x,r}$.
Thus it intertwines the one-sided normalized modes and all formal
coefficients of their half-vertices. Specified residue operations
require their own comparison below.

\begin{example}\label{v4-arith:curve:ex:curves}
For $C=\mathbb P^1$, $N_n=q^n+1$ and
$Z_C(t)=((1-t)(1-qt))^{-1}$.

For the smooth projective cubic
\[
 E/\mathbb F_2:\quad y^2+y=x^3+x+1,
\]
the first four point counts are $1,5,13,25$.
The equation $y^2+y=b$ has two solutions exactly when
$\operatorname{Tr}_{\mathbb F_{2^n}/\mathbb F_2}(b)=0$.
For $n=1,2,3,4$ there are respectively $0,2,6,12$ such values of
$x$, and the projective curve has one point at infinity. Consequently
\[
 Z_E(t)=1+t+3t^2+7t^3+15t^4+O(t^5).
\]

These counts determine the exchange through degree four. The full
formula
$Z_E(t)=(1-2t+2t^2)/((1-t)(1-2t))$
also uses the classical rationality and degree-two numerator theorem
for an elliptic curve.
\end{example}

\subsection{An elementary modification}

The full Hall algebra of coherent sheaves contains line bundles as
well as torsion sheaves. Fix a rational point $x$ and a line bundle
$L$. Extensions of $\mathcal O_x$ by $L$ have middle sheaf
$L\oplus\mathcal O_x$ or $L(x)$: the extension space is one
dimensional and all nonzero classes give $L(x)$.

In the split sheaf the relevant copies of $L$ are the graphs of the
$q$ maps $L\to\mathcal O_x$. In $L(x)$ there is the unique
subobject $L(x)(-x)=L$. Conversely
$\operatorname{Ext}^1(L,\mathcal O_x)=0$, and the split middle
sheaf has its unique torsion subobject. Hence
\[
 u_{\mathcal O_x}u_L=q\,u_{L\oplus\mathcal O_x}+u_{L(x)},
 \qquad u_Lu_{\mathcal O_x}=u_{L\oplus\mathcal O_x}.
\]

With $v=\sqrt q$, twist multiplication by
$u_A*u_B=v^{\chi(A,B)}u_Au_B$.
Here $\chi(\mathcal O_x,L)=-1$ and $\chi(L,\mathcal O_x)=1$, so
\[
 [u_{\mathcal O_x},u_L]_* =v^{-1}u_{L(x)}.
\]

On $\mathbb P^1$, summing over its $q+1$ rational points gives
\[
 \left[\sum_xu_{\mathcal O_x},u_{\mathcal O(n)}\right]_*
 =(v+v^{-1})u_{\mathcal O(n+1)}.
\]

This operation retains the elementary modification. A presentation
of a quantum loop algebra requires its other generators and relations.
An arithmetic trace comparison additionally requires its category,
state and functional.

\section{Local fields on every polynomial Hall state}
\label{v4-arith:curve:sec:local-fields}

Return to one point, with residue cardinality $Q$.
Set $h_{-r}=M_r$, $h_r=(Q^r-1)D_r=r\partial_{p_r}$ for $r>0$,
and $h_0=0$. Thus
\[
 [h_r,h_s]=r\delta_{r+s,0}I,
 \qquad h_r^{*G}=(Q^r-1)h_{-r}\quad(r>0).
\]

Define a derivation $T$ by $T(p_r)=rp_{r+1}$ and define
\[
 H(z)=\sum_{r\in\mathbb Z}h_rz^{-r-1}.
\]

For a polynomial monomial set
\begin{equation}\label{v4-arith:curve:eq:all-states}
 Y(p_{r_1}\cdots p_{r_l},z)
 =:\prod_{j=1}^l\frac{\partial_z^{r_j-1}H(z)}{(r_j-1)!}:,
 \qquad Y(1,z)=I.
\end{equation}
Normal ordering places creation modes to the left. This defines
the field on every state. Positive modes act locally finitely on
a polynomial input, and the creation series has a lower bound on
its powers of $z$.

\begin{proposition}\label{v4-arith:curve:prop:wick}
The fields \eqref{v4-arith:curve:eq:all-states}, with vacuum $1$
and translation $T$, give the free Heisenberg vertex algebra on
$H_Q$. Every finite ordered product has a common normally ordered
Wick expression, with rational contractions and finite poles along
pairwise diagonals. Every nesting expands that expression and
retains its nonconsecutive contractions.
\end{proposition}
\begin{proof}
Commuting an annihilator past a creator gives the scalar
$r\delta_{rs}$. Summing gives $(z-w)^{-2}$ in the expansion
$|w|<|z|$. For fixed polynomial states, commute each annihilator
through later creators. The terms are indexed by partial matchings
of the elementary factors. Matched factors contribute the
appropriate derivatives of $(z-w)^{-2}$ and unmatched factors
remain normally ordered. There are finitely many matchings and
derivative orders, so one rational denominator clears every term.

Regrouping partitions the same matchings according to the blocks
joined by their edges. Both nested expansions are therefore
expansions of the same rational expression. Empty matchings give
the vacuum identities.

Directly,
$[T,h_m]=-mh_{m-1}$, which proves translation covariance.
The two expansions of each contraction differ by a derivative of
the formal delta distribution. A finite power of $z-w$ kills the
commutator. These identities prove locality and associativity for
all the fields in \eqref{v4-arith:curve:eq:all-states}.
\end{proof}

Hall multiplication is represented by creation modes:
multiplication by $p_r$ is $h_{-r}$. It need not agree with the
vertex $(-1)$-product on composite states. For example,
\[
 (p_1^2)_{(-1)}p_1=p_1^3+2p_3,
 \qquad (p_1^2)p_1=p_1^3.
\]
The additional term is the constant coefficient of
$2z^{-2}\sum_{r>0}p_rz^{r-1}$ in $:H(z)^2:p_1$.

\subsection{Finite power covers}

At a closed point of degree $d$, pull back the current density
along $s=z^d$ on $\mathbb G_m$:
\[
 J_x(z)=dz^{d-1}H_x(z^d),\qquad
 J^x_n=\begin{cases}d\,h^x_{n/d},&d\mid n,\\0,&d\nmid n.
 \end{cases}
\]
The map is finite \'etale over $\mathbb C$ on the torus.
Its extension across zero is ramified, so the domain matters.
The common contraction is
\begin{equation}\label{v4-arith:curve:eq:cyclotomic}
 \frac{d^2z^{d-1}w^{d-1}}{(z^d-w^d)^2}
 =\sum_{\zeta^d=1}\frac{\zeta}{(z-\zeta w)^2}.
\end{equation}
Differentiate in $w$ the identity
$\sum_{\zeta^d=1}(z-\zeta w)^{-1}
=dz^{d-1}/(z^d-w^d)$ to obtain this formula.
Different point colors commute.

The coefficient of $z^{rd-1}$ in $J_x(z)1$ is $dp_{x,r}$.
Products of such creation coefficients recover all finite-support
Hall states. Normal products and derivatives have the pulled-back
Wick expansions. Derivative orders $a,b$ give pole order at most
$a+b+2$ at each divisor $z=\zeta w$. This is a bound for each
specified input. The states have no common finite derivative bound.

\subsection{Fields, quasi-modules and residue maps}

For a complex vector space $W$ write
$\mathcal E(W)=\operatorname{Hom}(W,W((z)))$.
A subset is \emph{quasi-local} if, for each pair of its fields,
a nonzero polynomial in their two coordinates annihilates the
commutator. It is \emph{quasi-compatible} if each finite ordered
product becomes a map to
$W[[z_1,\ldots,z_m]][z_1^{-1},\ldots,z_m^{-1}]$
after multiplication by a finite product of nonzero two-variable
polynomials. Formula~\eqref{v4-arith:curve:eq:cyclotomic} proves
quasi-locality of the point-labeled fields on $\mathcal F_C$.
Their finite Wick expressions prove quasi-compatibility.

Li's reconstruction theorem gives a vertex algebra \emph{of fields}
from a quasi-local subset of $\mathcal E(W)$, including its closure
under any chosen subgroup of dilations. Evaluation gives a faithful
quasi-module on $W$.\footnote{H.~Li, \emph{A new construction of
vertex algebras and quasi modules for vertex algebras}.
\href{https://arxiv.org/abs/math/0404110}{arXiv:math/0404110},
Theorems 4.9 and 5.4. The closure includes the identity field.
Faithfulness here follows from its inclusion in $\mathcal E(W)$.}

Concretely, if $p(z_1,z)$ clears the product, then
\begin{equation}\label{v4-arith:curve:eq:field-product}
 Y_{\mathcal E}(a(z),u)b(z)
 =\iota_{z,u}\bigl(p(z+u,z)^{-1}\bigr)
   \left.p(z_1,z)a(z_1)b(z)\right|_{z_1=z+u}.
\end{equation}

The symbol $\iota_{z,u}$ means expansion in
$\mathbb C((z))((u))$. Clearing first makes the substitution
defined. The theorem makes its coefficients independent of the
clearing polynomial and proves closure and the vertex identities.
Substitution at $z_1=\alpha z+u$ gives the dilation operations.
The original Hall space is the quasi-module, and has not been
identified with this vertex algebra of fields.

For finitely many point colors, the pole group is generated by
finitely many roots of unity. The whole point-labeled construction
is their filtered union: each operation uses finitely many
denominators and colors.

There are also actual supported residue maps on the power-cover
sector. Start with the Heisenberg fields of
Proposition~\ref{v4-arith:curve:prop:wick}. At the collision of two
positions, with $u=s_i-s_j$, the map on states and a multiplier $f$
is
\begin{equation}\label{v4-arith:curve:eq:residue}
 r_{ij}(f;\ldots,a,\ldots,b,\ldots)
 =\operatorname{pp}_u\bigl(fY(a,u)b\bigr).
\end{equation}

Here $\operatorname{pp}_u$ retains the negative powers of $u$.
The other states stay at their positions. Regular functions act
on principal parts by their finite normal Taylor expansions. For
example, for $t\ne0$ and $r,\ell\geq1$,
\[
 (u-t)^{-\ell}[u^{-r}]
 =(-1)^\ell\sum_{m=0}^{r-1}
   \binom{\ell+m-1}{m}t^{-\ell-m}[u^{m-r}].
\]
This retains the spectator poles and every finite normal derivative.

Translation covariance makes the maps horizontal for the state
connection $\partial_s+T$. Residues of disjoint pairs commute.
For three positions, the three iterated principal parts come from
the same rational Wick expression. The principal-parts relation
between its three collision channels is their residue identity.
Thus one has the full residue diagram, including its attachments,
before taking a cover. These statements concern the polynomial
Heisenberg states with their finite pole bounds.

Pull back this diagram along
\[
 (\mathbb G_m)^n\longrightarrow(\mathbb G_m)^n,\qquad
 (z_1,\ldots,z_n)\longmapsto(z_1^d,\ldots,z_n^d).
\]
Finite étale base change preserves its right differential modules,
residue maps and their relations. The diagonal $s_i=s_j$ pulls
back to the divisors $z_i=\zeta z_j$, with their actual
intersections and pulled-back attachments. Pullback of the
current density is exactly \eqref{v4-arith:curve:eq:cyclotomic}.
The cut tensor here excludes $z_i^d=z_j^d$ across a cut.
It would change if only $z_i=z_j$ were excluded.

For finitely many colors, use the independent local factors and
their labeled elementary insertions. Their finite Wick matchings
give the corresponding coefficient diagrams. Composite fields at
a retained root use \eqref{v4-arith:curve:eq:field-product}.
This supplies an operation comparison for this specified cover
sector. It leaves a field presentation of an arbitrary geometric
cut object as a separate question.

\section{Collective counting fields}
\label{v4-arith:curve:sec:counting}

Sum the point-labeled currents coefficientwise:
\[
 J(z)=\sum_xJ_x(z)=\sum_{n\in\mathbb Z}J_nz^{-n-1},
 \qquad J_{-n}=A_{-n},\quad J_n=(q^n-1)B_n\quad(n>0).
\]
Set $J_0=0$. A fixed mode involves only points whose degrees divide
its nonzero index, so it acts on the polynomial Hall space.
On a fixed input only finitely many positive modes act nontrivially.
The commutators are
\begin{equation}\label{v4-arith:curve:eq:counting-commutator}
 [J_n,J_{-m}]=\delta_{nm}nN_nI\qquad(n,m>0).
\end{equation}

\begin{proposition}\label{v4-arith:curve:prop:translation}
A linear operator $\mathsf T$ on this mode representation can
satisfy $[\mathsf T,J_n]=-nJ_{n-1}$ for all integers $n$ only if
$N_{n+1}=N_n$ for every $n\geq1$.
\end{proposition}
\begin{proof}
Apply $[\mathsf T,-]$ to $[J_{n+1},J_{-n}]=0$ and use Jacobi.
The resulting equality is
\[
 0=-(n+1)[J_n,J_{-n}]+n[J_{n+1},J_{-n-1}]
   =n(n+1)(N_{n+1}-N_n)I.
\]
The coefficient field has characteristic zero.
\end{proof}

For $\mathbb P^1$, the difference $N_{n+1}-N_n=q^n(q-1)$
is nonzero. Thus the unchanged collective current has no ordinary
vertex translation. The degree derivation
$L_0(p_{x,r})=rd_xp_{x,r}$ does satisfy
$[L_0,J(z)]=(z\partial_z+1)J(z)$.
It implements dilation.

\subsection{Reciprocal normalization}

Write a rational counting function with constant term one as
\[
 Z(t)=\prod_{\lambda\in\Lambda}(1-\lambda t)^{-\epsilon_\lambda},
 \qquad \epsilon_\lambda\in\mathbb Z,
\]
with finitely many nonzero parameters. Its current contraction is
\begin{equation}\label{v4-arith:curve:eq:counting-kernel}
 C(z,w)=\frac1{zw}(t\partial_t)^2\log Z(t)
       =\sum_{\lambda\in\Lambda}
           \frac{\epsilon_\lambda\lambda}{(z-\lambda w)^2},
 \qquad t=w/z.
\end{equation}
This follows by twice differentiating each
$-\log(1-\lambda t)$.

For the curve, the classical rationality theorem gives
\[
 Z_C(t)=\frac{\prod_{i=1}^{2g}(1-\alpha_it)}{(1-t)(1-qt)}.
\]

The numerator parameters $\alpha_i$ occur in the current
contraction as well as $1,q$. They must therefore occur among
its allowed collision gains. Poles of the current contraction
are not determined just by poles of $Z_C$.

The full Laurent field needs both orders of a contraction.
The expression \eqref{v4-arith:curve:eq:counting-kernel} need not
be symmetric under $z,w$ interchange, so quasi-locality of each
$J_x$ alone does not establish quasi-locality of the infinite sum.
Rationality and reciprocal eigenvalue pairing give a different
collective field.

Choose $\sqrt q>0$ and define
\[
 \widetilde J_{-n}=A_{-n},\qquad
 \widetilde J_n=q^{-n/2}(q^n-1)B_n,\qquad
 c_n=q^{-n/2}N_n.
\]
These formulas use $n>0$. Set $\widetilde J_0=0$.

The signed multiset
\[
 \widetilde\Lambda
   =\{q^{-1/2},q^{1/2};\ \alpha_1/q^{1/2},\ldots,\alpha_{2g}/q^{1/2}\}
\]
has signs $+,+,-,\ldots,-$. The pairing
$\alpha\leftrightarrow q/\alpha$ makes this multiset invariant
under inversion. This uses reciprocal duality. No assertion about
the absolute values of the $\alpha_i$ is needed.

\begin{theorem}\label{v4-arith:curve:thm:reciprocal}
The field $\widetilde J(z)=\sum_n\widetilde J_nz^{-n-1}$
has the symmetric rational contraction
\[
 \widetilde C(z,w)=
 \sum_{\lambda\in\widetilde\Lambda}
       \frac{\epsilon_\lambda\lambda}{(z-\lambda w)^2}.
\]
Its commutator is annihilated by
$\prod_{\lambda\in\widetilde\Lambda}(z-\lambda w)^2$.
Its dilates generate a quasi-local system, and hence a vertex
algebra of fields with faithful quasi-module $\mathcal F_C$.
\end{theorem}
\begin{proof}
Inversion invariance and
\[
 \frac{\lambda}{(w-\lambda z)^2}
 =\frac{\lambda^{-1}}{(z-\lambda^{-1}w)^2}
\]
give symmetry. The coefficient of
$z^{-n-1}w^{n-1}$ in the forward expansion is $nc_n$.
Inversion invariance gives the same constants in the reverse
expansion.

Their difference is
\[
 \sum_{\lambda}\epsilon_\lambda\partial_w
       \bigl(z^{-1}\delta(\lambda w/z)\bigr),
 \qquad \delta(t)=\sum_{m\in\mathbb Z}t^m.
\]
Each summand is killed by $(z-\lambda w)^2$.
Every finite set of dilates has a finite set of these factors.
The quasi-local reconstruction theorem therefore applies.
\end{proof}

Translation in this vertex algebra acts on its fields through
\eqref{v4-arith:curve:eq:field-product}. The Hall space is a
quasi-module, and its new positive modes also differ from $J_n$.
Thus the theorem imposes no impossible translation on the
unchanged representation in
Proposition~\ref{v4-arith:curve:prop:translation}.
The collective modes retain the counting sequence. Individual
closed-point labels remain in the separate fields $J_x$.

The point-labeled and aggregate rational fields together are
quasi-compatible. For a fixed point of degree $d$, their mixed
contraction uses only $n\in d\mathbb Z$. Its coefficients are
finite combinations of $n$ times geometric progressions.
It therefore sums to a rational function. Aggregate contractions
use \eqref{v4-arith:curve:eq:counting-kernel} and its rescalings.

Wick expansion clears a finite common denominator in every finite
ordered product. Li's nonlocal reconstruction then gives a
nonlocal vertex algebra of fields and a faithful Hall
quasi-module.\footnote{H.~Li, \emph{Nonlocal vertex algebras
generated by formal vertex operators}.
\href{https://arxiv.org/abs/math/0502244v2}{arXiv:math/0502244v2},
Theorem 2.20. Quasi-compatibility is Definition 2.10.}

Pairwise quasi-locality of this combined set has not been assumed.
The native Green current has an infinite-product kernel.
It uses the analytic recipient below.

\subsection{A finite-dimensional trace comparison}

For a finite-dimensional graded vector space $H=\bigoplus_iH^i$
and a grading-preserving endomorphism $F$, graded symmetry gives
\[
 \sum_{n\geq0}\operatorname{str}(F;\operatorname{Sym}^nH)t^n
 =\prod_i\det(1-tF|H^i)^{(-1)^{i+1}}.
\]
Indeed, after extending scalars and triangularizing, an even
eigenvector contributes $(1-\alpha t)^{-1}$. An odd eigenvector
generates an exterior algebra and contributes $1-\alpha t$.
Multiply the factors and descend the coefficients. Jordan
off-diagonal entries do not change the traces.

The two even eigenvalues $1,q$, in degrees zero and two, reproduce
the series of $\mathbb P^1$. A more discriminating example is
$y^2=x^3-x$ over $\mathbb F_5$, which has seven affine points
and one point at infinity. The matrices
\[
 F=\begin{pmatrix}0&-5\\1&-2\end{pmatrix},
 \qquad H=\begin{pmatrix}1&-1\\-1&5\end{pmatrix}
\]
satisfy $F^tHF=5H$. The leading principal minors of $H$ are
$1,4$, so $H$ is positive definite. This particular operator
has eigenvalues of modulus $\sqrt5$ and determinant polynomial
$1+2t+5t^2$.

Identifying it with a geometric Frobenius
representation still requires the cohomological realization.
The point count and positive matrix alone do not construct it.

\section{Gain equations and their derived intersections}
\label{v4-arith:curve:sec:gains}

Let $k$ have characteristic zero and let $G\subset k^\times$ be
a subgroup containing the chosen collision gains. Work on
$X=(\mathbb G_m)^n$, with coordinate ring
$S=k[z_1^{\pm1},\ldots,z_n^{\pm1}]$.
For a finite oriented graph $\Gamma$, including repeated edges,
an edge $e=(i,j,\lambda_e)$ names the equation
$f_e=z_i-\lambda_ez_j$. A component is \emph{balanced} if the
oriented product of gains around every cycle is one. Reversing
an edge replaces its gain by its inverse.

The scalar derived intersection is the Koszul complex
\begin{equation}\label{v4-arith:curve:eq:gain-complex}
 K_\Gamma=(S\otimes\Lambda(\xi_e:e\in E(\Gamma)),d),
 \qquad |\xi_e|=-1,\qquad d\xi_e=f_e.
\end{equation}
For pole coefficients one instead inverts a finite collection
of nonzero $f_e$. The union of these localization rings records
functions with finitely many allowed poles. It does not assert
that an infinite arrangement is a finite-type scheme.

\begin{theorem}\label{v4-arith:curve:thm:gains}
Suppose $\Gamma$ has $n$ vertices, $e$ edges and $c$ components.
If all components are balanced, a spanning forest gives
\[
 K_\Gamma\simeq
 k[y_1^{\pm1},\ldots,y_c^{\pm1}]
     \otimes\Lambda(\eta_1,\ldots,\eta_{e-n+c}),
 \qquad |\eta_j|=-1,\quad d\eta_j=0.
\]
If a component is unbalanced, $K_\Gamma$ is contractible.
\end{theorem}
\begin{proof}
Removing leaves of a spanning forest eliminates one coordinate
at each step by a nonzero scalar multiple of its neighbor.
Its $n-c$ equations form a regular sequence, with quotient the
torus on the root coordinates. In a balanced component, every
remaining equation is a constant linear combination of forest
equations. Subtract this combination from its exterior generator.
The resulting invertible triangular change splits off one closed
generator for each extra edge.

An unbalanced cycle instead leaves $(1-\lambda)y$ with
$\lambda\ne1$ and $y$ invertible. Thus its equations generate
the unit ideal. If $\sum_eg_ef_e=1$, then exterior multiplication
by $\beta=\sum_eg_e\xi_e$ satisfies $d(\beta\wedge-)+\beta\wedge d=1$.
This contracts the complete complex.
\end{proof}

For the balanced triangle,
\begin{equation}\label{v4-arith:curve:eq:gain-triangle}
 \begin{gathered}
 f_1=z_1-\lambda z_2,\quad f_2=z_2-\mu z_3,\quad
 f_3=z_1-\lambda\mu z_3=f_1+\lambda f_2,\\
 \eta=\xi_3-\xi_1-\lambda\xi_2,\qquad d\eta=0.
 \end{gathered}
\end{equation}
Quotient by the forest equations and send $\xi_1,\xi_2$ to zero.
The image of $\eta$ is a nonzero exterior generator, so it cannot
have been a boundary. If the third gain is instead $\nu\ne\lambda\mu$,
the explicit primitive of the unit is
\[
 \beta=\frac{\xi_3-\xi_1-\lambda\xi_2}{(\lambda\mu-\nu)z_3},
 \qquad d\beta=1.
\]

The forest basis computes the cohomology. Equation operations
belong to the unreduced complex. A ring map $S\to S'$ and a
matrix $U=(u_{te})$ satisfying
$f_e\mapsto\sum_tu_{te}f'_t$ give the cochain map
$\xi_e\mapsto\sum_tu_{te}\xi'_t$. Composition is matrix
multiplication with coefficient substitution.

For a clustering $\alpha:I\twoheadrightarrow J$ and scales
$\gamma_i\in G$, the substitution $z_i=\gamma_i y_{\alpha(i)}$ gives
\[
 f_{ij,\lambda}\longmapsto
 \gamma_i\bigl(y_{\alpha(i)}
       -\lambda\gamma_j\gamma_i^{-1}y_{\alpha(j)}\bigr).
\]

The exterior generator has the same factor $\gamma_i$.
An internal zero equation retains its exterior generator.
An internal unit equation contracts. An equation that becomes
zero cannot remain inverted, so a collision requires its
restriction complex.

\subsection{The divided exponential chart}

In this spectral coefficient family, regard $q$ as a formal variable.
For gains $q^a$, with integers $a$, introduce a formal parameter
$\hbar$ and use $q=e^\hbar$, $z_i=e^{\hbar x_i}$.
In $k[x_1,\ldots,x_n][[\hbar]]$ the divided equation is
\[
 \frac{e^{\hbar x_i}-e^{\hbar(x_j+a)}}{\hbar}
 =(x_i-x_j-a)U_{ij,a}(\hbar),\qquad U_{ij,a}(0)=1.
\]
The quotient denotes a power series, with no inversion of $\hbar$.

Rescaling the exterior generator by the unit $U_{ij,a}^{-1}$
gives the affine equation. A forest quotient is
$k[y_1,\ldots,y_c][[\hbar]]$, which is flat over $k[[\hbar]]$.
Balanced cycles retain the excess exterior factors. An unbalanced
cycle telescopes to a nonzero integer, so its complex contracts.
After inverting $\hbar$, this is the pulled-back multiplicative
arrangement up to units.

The special fiber is the shifted arrangement $x_i-x_j=a$.
Undivided equations would all reduce to zero. This constructs a
family of spectral coefficient arrangements. Extending Green
operators across $\hbar=0$ further requires a lattice and a
normalization, since $q^n-1$ becomes $\hbar$ times a unit.
No family of arithmetic curves is supplied by the coordinate
substitution.

\subsection{Scalar operations on the pole ring}

Fix colors $\alpha$ and rational kernels
$K_{\alpha\beta}(z,w)$ with finitely many allowed gain poles.
For functions on two ordered lists of colors define
\[
 (f\diamond g)(z_1,\ldots,z_{m+n})
 =f(z_1,\ldots,z_m)g(z_{m+1},\ldots,z_{m+n})
   \prod_{i\leq m<j}K_{\alpha_i\alpha_j}(z_i,z_j).
\]
This is an associative product, with the empty list as unit:
either bracketing of three blocks includes each cross-block
factor exactly once. Summing over order-preserving shuffles
gives an associative product on symmetric colored functions.

For graded colors, include the Koszul permutation sign.
Both bracketings have the same three-block shuffles and signs.

Inverse exchanges may require gains from zeros as well as poles
of the original kernels. The scalar construction then uses the
enlarged group. A Hall comparison must also supply the maps of
coefficients, orientations and any numerator conditions required
by its recipient. The scalar product proves none of these from
an equality of kernels.

\section{Supported coefficients with excess}
\label{v4-arith:curve:sec:supported}

For one regular equation $f$, the scalar Koszul complex and the
localization complex have different purposes. The latter is
$L_f=[S\to S[f^{-1}]]$ in degrees $0,1$.
Its map is horizontal for the usual left differential action.
Localization gives
\[
 L_f\otimes_S^{\mathbb L}L_f\simeq L_f:
\]
tensor the triangle $L_f\to S\to S[f^{-1}]$ with $L_f$ and use
the acyclicity of $L_f\otimes_SS[f^{-1}]$.
In contrast, $S/(f)\otimes_S^{\mathbb L}S/(f)$ has a nonzero
degree-minus-one Tor class.

For $f=z_i-\lambda z_j$, $i\ne j$, the Koszul complex
$[S\xrightarrow{f}S]$ in degrees $-1,0$ cannot admit a
differential-module enhancement with the same object of $D(S)$.
Such an enhancement would give its cohomology $S/(f)$ the stated
$S$-action and a left differential action. But $f$ acts by zero,
whereas $[\partial_{z_i},f]=1$ would force their commutator to
be the identity on this nonzero quotient.

We construct a supported recipient that also retains the excess.
For a balanced graph let $V=k^{E(\Gamma)}$, with basis $\xi_e$,
and view the equations as linear forms on $A_0=k^n$. Put
\begin{equation}\label{v4-arith:curve:eq:linear-data}
 F:V\longrightarrow A_0^*,\quad F(\xi_e)=f_e,\qquad
 W=\operatorname{im}F,\quad E=\ker F,\quad T_\Gamma=W^\perp .
\end{equation}

Here $E$ denotes the space of equation relations. If $r=\dim W$
and $c=\dim T_\Gamma$, then $r=n-c$ and
$Z=T_\Gamma\cap X\cong(\mathbb G_m)^c$.

Let $L=\det T_\Gamma^*$ be its constant tangent density line.
The density-valued scalar source is
$\mathcal K_\Gamma=L\otimes K_\Gamma$, with
\[
 H^{-p}(\mathcal K_\Gamma)=\omega_Z\otimes_k\Lambda^pE,
 \qquad \omega_Z=\mathcal O_Z\otimes_k L.
\]

These cohomology identifications are independent of a forest:
wedge closed vectors of $E$ with a degree-zero class. A forest
splitting proves that this map is an isomorphism.

All complexes are cohomological. $C[m]^j=C^{j+m}$ with
differential $(-1)^m d_C$.

\subsection{Normal derivatives and the ordinary module}

For $i:Z\hookrightarrow X$ define the right differential module
\[
 \mathcal N=i_+\omega_Z
 =i_*\bigl(\omega_Z\otimes_{\mathcal D_Z}\mathcal D_{Z\to X}\bigr),
 \qquad
 \mathcal D_{Z\to X}
 =\mathcal O_Z\otimes_{i^{-1}\mathcal O_X}i^{-1}\mathcal D_X.
\]
It lies in cohomological degree zero.

Choose complementary linear
coordinates $(u_1,\ldots,u_r,y_1,\ldots,y_c)$ with $Z=(u=0)$,
and set $\tau=dy_1\wedge\cdots\wedge dy_c$.
On the actual open chart its sections have the normal form
\begin{equation}\label{v4-arith:curve:eq:normal-jets}
 \mathcal N=\bigoplus_{\alpha\in\mathbb N^r}
       \mathcal O_Z\,\delta_{\tau,\alpha},\qquad
 \delta_{\tau,\alpha}
 =\left[\frac{du_1\wedge\cdots\wedge du_r\wedge\tau}
                    {\prod_j u_j^{\alpha_j+1}}\right].
\end{equation}

The brackets denote local cohomology in degree $r$, placed in
degree zero in $\mathcal N$. The residue of
$du_1\wedge\cdots\wedge du_r/(u_1\cdots u_r)$ is one.
Equivalently,
$\delta_{\tau,\alpha}=\delta_{\tau,0}\partial_u^\alpha/\alpha!$.
Thus
\begin{equation}\label{v4-arith:curve:eq:jet-actions}
 \begin{aligned}
 \delta_{\tau,\alpha}u_j&=\delta_{\tau,\alpha-e_j},&
 \delta_{\tau,\alpha}\partial_{u_j}
       &=(\alpha_j+1)\delta_{\tau,\alpha+e_j},\\
 (b\delta_{\tau,\alpha})\partial_{y_j}
       &=-\partial_{y_j}(b)\delta_{\tau,\alpha}.
 \end{aligned}
\end{equation}
The first expression is zero if $\alpha_j=0$.

These formulas follow from
$(a\Omega)\partial=-\partial(a)\Omega$ for a constant top form
$\Omega$. They satisfy the right Weyl relation
$(m\partial_{u_i})u_j-(mu_j)\partial_{u_i}=\delta_{ij}m$.
Every section has finite normal order. All finite orders occur.
No infinite normal-derivative tail is introduced.

A constant change $u'=Au$ preserves the normalized order-zero
residue and transforms derivatives by
$\partial_{u'_i}=\sum_j(A^{-1})_{ji}\partial_{u_j}$.
Both numerator density and denominator must transform.
A tangent change $y'=By$ multiplies the density by $\det B$.
These rules follow by testing the fractions against normal
polynomials: $\delta_{\tau,\alpha}$ extracts the coefficient
of $u^\alpha$ with its tangent density.

If $I=(f_e)$, the actual support complex is
\[
 C_I(\omega_X)=
 \bigotimes_{e\in E(\Gamma)}[S\longrightarrow S[f_e^{-1}]]
       \otimes_S\omega_X.
\]
It has only $H^r=\mathcal N$. For forest equations this is the
tensor product of the one-variable principal-parts calculation.

A redundant $f\in I$ adds no cohomology.
Localization of the supported finite jets at $f$ is zero.
Thus tensoring with $[S\to S[f^{-1}]]$ changes the complex
by an acyclic cone.
For a zero equation, the localization is zero and the extra
factor is just $S$ in degree zero. Therefore
$C_I(\omega_X)[r]\simeq\mathcal N$, including repeated equations.

There is a scalar cochain map
\[
 \varepsilon:\mathcal K_\Gamma\longrightarrow\operatorname{Res}_S\mathcal N,
 \qquad \tau a\longmapsto a|_Z\delta_{\tau,0},
\]
which is zero on positive exterior length.
Normal annihilation makes it $S$-linear and a cochain map.
It includes $\omega_Z$ as the normal socle, the part killed by $I$.

Every degree-zero morphism from $\mathcal K_\Gamma$ to this
ordinary target kills negative cohomology. More generally, $\mathcal N[m]$ has cohomology only in degree $-m$.
It cannot receive both a nonzero vacuum class in degree zero
and a nonzero excess class in degree $-1$.
Replacing complexes by resolutions, changing density by a unit,
or retaining more normal derivatives leaves this degree argument
unchanged.

\subsection{A derived supported complex}

Let $D=\Gamma(X,\mathcal D_X)$. Define
\begin{equation}\label{v4-arith:curve:eq:spencer}
 \mathcal P_\Gamma=L\otimes_k\Lambda(V\oplus T_\Gamma)\otimes_kD,
 \qquad |V|=|T_\Gamma|=-1.
\end{equation}
For $v\in V$ and $t\in T_\Gamma$ put
$\mathsf a(v)=F(v)$ and $\mathsf a(t)=\partial_t$.
On an ordered wedge the differential is
\[
 d(\tau\otimes a_1\wedge\cdots\wedge a_m\otimes Q)
 =\sum_{j=1}^m(-1)^{j-1}\tau\otimes
    a_1\wedge\cdots\widehat a_j\cdots\wedge a_m
          \otimes\mathsf a(a_j)Q.
\]
Multiplication by $\mathsf a(a_j)$ is on the left of $Q$.
The differential is right $D$-linear.

The operators commute because
$[\partial_t,F(v)]=F(v)(t)=0$, so $d^2=0$.
The tangent exterior factors impose the Spencer relations.

There is an actual $S$-linear cochain map
\begin{equation}\label{v4-arith:curve:eq:spencer-map}
 j_\Gamma:\mathcal K_\Gamma\longrightarrow\operatorname{Res}_S\mathcal P_\Gamma,
 \qquad
 \tau a\,v_1\wedge\cdots\wedge v_p
 \longmapsto\tau\otimes(v_1\wedge\cdots\wedge v_p)\otimes a .
\end{equation}
Functions commute with the $F(v_i)$, which proves its cochain
equation. The scalar source is not being given a differential
action by this map.

\begin{theorem}\label{v4-arith:curve:thm:spencer}
There are canonical right $D$-module isomorphisms
\[
 H^{-p}(\mathcal P_\Gamma)\cong\Lambda^pE\otimes_k\mathcal N.
\]
On cohomology, $j_\Gamma$ includes
$\omega_Z\otimes\Lambda^pE$ at normal order zero. This inclusion
is injective and its right $D$-span is the whole group.
It retains the vacuum, every excess class and, by differentiation,
every finite normal descendant.
\end{theorem}
\begin{proof}
Choose $V_0$ complementary to $E$ and a basis of $T_\Gamma$.
The closed factor $\Lambda E$ splits off.
The remaining right Koszul complex uses the commuting operators
$u_1,\ldots,u_r,\partial_{y_1},\ldots,\partial_{y_c}$.

Filter $D$ by differential order, assigning weight zero to
normal exterior generators and one to tangent exterior generators.
Its symbol complex is the commutative Koszul complex of
\[
 u_1,\ldots,u_r,p_{y_1},\ldots,p_{y_c}
 \quad\hbox{in }S[p_1,\ldots,p_n].
\]
This is a regular sequence. To solve a negative-degree cycle,
lift a boundary for its leading symbol and subtract it.
The filtration degree decreases and is bounded below in every
cochain degree, so the process terminates.

The degree-zero quotient is
\[
 L\otimes_kD\big/\left(\sum_i u_iD+\sum_j\partial_{y_j}D\right).
\]
Normal ordering identifies it with the transfer module
$i_+\omega_Z$: the tangent density relations remove tangent
derivatives and the normal derivatives give
\eqref{v4-arith:curve:eq:normal-jets}.
Wedging with closed elements of $E$ gives a canonical map
$\Lambda^pE\otimes H^0(\mathcal P_\Gamma)\to H^{-p}(\mathcal P_\Gamma)$.
The splitting proves that it is an isomorphism, removing the
choice of complement from the conclusion.
Formula~\eqref{v4-arith:curve:eq:spencer-map} sends the source to
the normal socle. The jet basis proves injectivity, and right
normal differentiation generates the entire target cohomology.
\end{proof}

Projection in exterior degree zero defines a right $D$-linear
map $\pi:\mathcal P_\Gamma\to\mathcal N$ by
$\pi(\tau\otimes Q)=\delta_{\tau,0}Q$.
Its cochain equation is exactly the normal and tangent
annihilation of $\delta_{\tau,0}$.
Thus $\pi j_\Gamma=\varepsilon$.
If $E\ne0$, $\pi$ is not a quasi-isomorphism: it kills the
negative excess.

Universal differential induction is a further useful recipient,
\[
 \mathcal U_\Gamma=\mathcal K_\Gamma\otimes_SD
                 =L\otimes\Lambda V\otimes D.
\]
Normal ordering makes $D$ free as a left $S$-module. Induction
therefore preserves the scalar cohomology and is functorial for
every $S$-linear equation map, including matrices with variable
coefficients and unchanged density. Inclusion of tangent exterior
degree zero gives $\mathcal U_\Gamma\to\mathcal P_\Gamma$.
The latter complex imposes the tangent relations while retaining
the excess.

In one normal coordinate, let $A=k[u]$ and
$E_u=A[u^{-1}]/A$. Left-to-right conversion gives
$(u^{-n}\otimes du)\partial_u=n u^{-n-1}\otimes du$.
Thus $\delta_m=u^{-m-1}\otimes du$ agrees with
\eqref{v4-arith:curve:eq:jet-actions}, and
$[D\xrightarrow{u\cdot}D]$ in degrees $-1,0$ is its free
right $D$-resolution. This compares the supported normal
coefficient resolutions without changing their derivative actions.

\subsection{Supported dual functionals}

The lost excess still has a supported dual functional. On the
linear normal chart, with $B=\mathcal O_Z$, one has
\begin{equation}\label{v4-arith:curve:eq:supported-dual}
 H^p\operatorname{RHom}_S(\mathcal K_\Gamma,\mathcal N)
       \cong B\otimes_k\Lambda^pE^*,\qquad 0\leq p\leq\dim E.
\end{equation}
All other degrees vanish.

To compute it, the bounded free source
allows ordinary Hom. Its degree-$p$ differential is
$-(-1)^p\phi d$. Rescaling the dual basis in degree $p$ by
$(-1)^{p(p+1)/2}$ gives the normal cochain Koszul differential.

For one normal variable this is
$[\mathcal N\xrightarrow{u}\mathcal N]$ in degrees $0,1$.
Multiplication by $u$ is surjective, with kernel $B\delta_0$.
The right inverse $R\delta_m=\delta_{m+1}$ satisfies
$uR=1$ and $Ru=1-\iota p$, where $p$ extracts $\delta_0$.

These homotopies are $B$-linear. The socle inclusion is
$S$-linear and a quasi-isomorphism. Tensoring the normal
contractions, with their cochain signs, leaves the socle and
$\Lambda E^*$. The density line cancels, proving
\eqref{v4-arith:curve:eq:supported-dual}.

For the triangle the functional
\[
 \ell(dy\,\xi_3)=\delta_{dy,0},\qquad
 \ell(dy\,\xi_1)=\ell(dy\,\xi_2)=0
\]
is closed and evaluates nontrivially on $dy\,\eta$.
A Hom boundary evaluates to zero on that closed class.
It therefore defines a nonzero map
$\mathcal K_\Gamma\to\mathcal N[1]$, which kills the
degree-zero vacuum. Derivative order and cohomological degree
remain different gradings.

Keeping differential-module maps adds the tangent de Rham
complex. On $Z=(\mathbb G_m)^c$ the bounded free complex
\eqref{v4-arith:curve:eq:spencer} gives
\[
 \operatorname{RHom}_D(\mathcal P_\Gamma,\mathcal N)
 \simeq(\Omega^\bullet(Z),d_{\mathrm{dR}})\otimes_k\Lambda(E^*),
 \qquad |E^*|=1.
\]

The same normal contractions commute with tangent
differentiation. The tangent Spencer dual has differential
$-d_{\mathrm{dR}}$ after the preceding degree normalization.
Multiplying tangent degree $a$ by $(-1)^a$ gives $d_{\mathrm{dR}}$.
The resulting degree-$m$ cohomology is
$\bigoplus_{a+p=m}H^a_{\mathrm{dR}}(Z)\otimes\Lambda^pE^*$.

Indeed $d(y^j)=j y^j d\log y$ contracts every nonzero Laurent
weight. Weight zero leaves $1,d\log y$, and tensoring gives
$H^\bullet_{\mathrm{dR}}(Z)=\Lambda^\bullet\langle d\log y_1,\ldots,d\log y_c\rangle$.

\subsection{Equation operations and tangent connections}

Suppose $F:V\to A_0^*$ and $F':V'\to A_0^*$ have the same image,
and let $U:V\to V'$ be constant with $F=F'U$. The assignments
\[
 v\longmapsto Uv,\qquad t\longmapsto t,\qquad
 Q\longmapsto Q,\qquad \tau\longmapsto\tau
\]
define a right $D$-linear map $\mathcal P_F\to\mathcal P_{F'}$.
The normal differentials agree by $F=F'U$, and constant
coefficients commute with tangent differentiation. Thus the map
is a cochain map and the square with the two maps $j$ commutes
strictly. It acts on cohomology as
$\Lambda^p(U|_E)\otimes I_{\mathcal N}$, so it intertwines the
excess action and composes by matrix multiplication.
It may kill an excess class when $U|_E$ does.

An invertible linear coordinate map gives the same comparison
on open domains that it actually identifies. It pulls back the
density, sends tangent vectors by the inverse coordinate map,
and transforms every normal derivative by the chain rule.
These operations compose strictly. Diagonal scalings and
coordinate permutations preserve the torus. An arbitrary
invertible linear map need not preserve it.

For example, in the triangle take $\lambda=2,\mu=4$ and
$u=z_1-2z_2$, $v=z_2-4z_3$, $y=z_3$. Under
\[
 z_1=4w_1,\qquad z_2=\tfrac12 w_2,\qquad z_3=w_3,
\]
put $u'=w_1-\tfrac14w_2$, $v'=w_2-8w_3$, $y'=w_3$.
The equations give $u=4u'$, $v=\tfrac12v'$ and
$\eta\mapsto4\eta'$, where
$\eta'=\xi'_3-\xi'_1-\tfrac14\xi'_2$.

The density-normalized jets transform by
\[
 \delta_{dy,(m,n)}\longmapsto4^{-m}2^n\delta_{dy',(m,n)}.
\]
Indeed the normal Jacobian is $2$ and the denominator factor
is $4^{m+1}(1/2)^{n+1}$. The vacuum is unchanged. Omitting
the Jacobian would already multiply it by $1/2$.

The tangent connection restricts variable equation operations.
On the same triangle let an equation automorphism fix
$\xi_1,\xi_2$ and send
\[
 \xi_3\longmapsto(1-y)\xi_1+\lambda(1-y)\xi_2+y\xi_3.
\]
It fixes the equations and sends $\eta$ to $y\eta$.

In supported cohomology, however,
$m=\eta\delta_{dy,0}$ satisfies $m\partial_y=0$ whereas
$(ym)\partial_y=-m\ne0$. No right $D$-linear cochain map on
the fixed $\mathcal P_\Gamma$ can induce this action.
Homotopies do not change the induced cohomology map.

Raw induction $\mathcal U_\Gamma$ remains functorial for this
scalar equation map. Removing its free tangent derivatives
while retaining the map requires additional connection data.

Equation symmetries can give a second obstruction, even when
determinants are unchanged. Repeat the third triangle edge and
put
$\eta_i=\xi_{i+2}-\xi_1-\lambda\xi_2$, $i=1,2$.
Fix the forest generators and act by
$\eta_1\mapsto t\eta_1$, $\eta_2\mapsto t^{-1}\eta_2$.
This automorphism fixes the equations, ambient coordinates,
normal quotient and tangent density, and has determinant one.

If a target assigns it the identity at $t=2$, a natural
homogeneous comparison kills both excess classes.
Their images would have to equal twice and half themselves.
Equation, tangent and normal determinant twists leave this
argument unchanged. The invariant top wedge is not killed
by this argument.

For the ordinary module on the split torus,
$\operatorname{End}_D(\mathcal N)=k$.
An endomorphism sends its normal socle generator to
$b\delta_{\tau,0}$. The normal relations force the socle,
and the tangent relations force every derivative of $b$ to vanish.
Right normal differentiation then determines the whole map.

Consequently a coherent action of the full $\operatorname{SL}(E)$
for $\dim E=2$ on any single shift $\mathcal N[m]$ is a character
to $k^\times$, and is trivial. To see this last assertion directly,
conjugation by $\operatorname{diag}(2,1/2)$ makes every upper
or lower elementary unipotent a commutator. These unipotents
generate $\operatorname{SL}_2(k)$.

Thus this naturality obstruction applies to the ordinary target
with its full equation action. A derived target carrying
$\Lambda E$ has the nontrivial excess action instead.

With three copies of the third edge, a cyclic permutation is even.
A natural comparison with identity target action kills the
nonzero difference $\xi_3-\xi_4=(1-\sigma)\eta_1$.

\subsection{A nontransverse clustering}

Again take the gains $2,4,8$ and restrict to
\[
 \iota_Y:Y=(\mathbb G_m)^2\longrightarrow X,\qquad
 (x,y)\longmapsto(2x,x,y).
\]
This is $u=0$, and $g=x-4y$ is the surviving normal coordinate.
The three equations become $(0,g,2g)$.
Writing $\zeta=\xi_1$ and
$\eta=\xi_3-\xi_1-2\xi_2$, one gets
\[
 \iota_Y^*K_\Gamma=K(g)\otimes\Lambda(\zeta,\eta).
\]
The excess $\eta$ is inherited. $\zeta$ is created by this
nontransverse restriction.

Let $\mathcal N_X$ have normal coordinates $(u,v)$ and let
$\mathcal N_Y$ have normal coordinate $g=v$, both with density
$dy$ along the common stratum. For the derived pullback of
underlying modules,
\begin{equation}\label{v4-arith:curve:eq:cluster-pullback}
 \mathcal O_Y\otimes_{\mathcal O_X}^{\mathbb L}
          \operatorname{Res}\mathcal N_X
 \simeq\mathcal N_Y[1].
\end{equation}
This convention is an $\mathcal O$-module pullback, rather than
a dimension-normalized differential inverse image.
Resolve $\mathcal O_Y$ by
$[\mathcal O_X\xrightarrow{u}\mathcal O_X]$.
After tensoring, multiplication by $u$ is surjective on
$\mathcal N_X$ and has kernel
$\bigoplus_{n\geq0}k[y^{\pm1}]\delta_{dy,(0,n)}$.
The residual $v,\partial_v,y,\partial_y$ actions are exactly
those of $\mathcal N_Y$, proving
\eqref{v4-arith:curve:eq:cluster-pullback} and retaining every
finite $v$-descendant.

Under the derived pullback of $\varepsilon$, the new class
$dy\,\zeta$ maps to minus the supported vacuum in degree $-1$.
Indeed, if $d\theta=u$ is the resolution generator, then
$dy(\xi_1-\theta)$ is a closed lift and maps to
$-\theta\delta_{dy,0}$. The inherited $dy\,\eta$ already maps
to zero before restriction, so still maps to zero.
Thus inherited excess and a new Tor class have different
images. No evaluation of a distribution on its singular
support has been used.

Changing the third gain to $16$ gives
$d(\xi_3-\xi_1-2\xi_2)=-8z_3$.
The primitive $-(\xi_3-\xi_1-2\xi_2)/(8z_3)$ of the unit
remains a primitive after this clustering, with $z_3=y$.
Both scalar complexes contract and both supports are empty.

\subsection{Contact differentials and the geometric comparison}

A meromorphic coefficient map also needs its contact
differential. Let $C,\mathcal L,J$ be cochain complexes,
let $W:C\to\mathcal L$ be a cochain map, and suppose a supported
extension has graded form $\mathcal L\oplus J$ and differential
\[
 d(a,s)=(d_{\mathcal L}a,d_Js+Ra),\qquad
 d_JR+Rd_{\mathcal L}=0.
\]
Its pullback along $W$ has the actual differential
\[
 E_W=C\oplus J,\qquad d(c,s)=(d_Cc,d_Js+RWc).
\]
The map $(c,s)\mapsto(Wc,s)$ is a cochain map because
$d_JRW+RWd_C=0$. A map to the localized component alone
omits this term.

With two independent normals, the localization
complex has both faces and their corner:
\[
 S\longrightarrow S[f^{-1}]\oplus S[g^{-1}]
       \longrightarrow S[(fg)^{-1}],
 \qquad d_0(a)=(a,a),\quad d_1(b,c)=c-b.
\]
The sign makes $d_1d_0=0$ and determines the corresponding
connecting signs in the supported extension.

The maps $\varepsilon$, $j_\Gamma$ and the supported dual
functional now have different, explicit meanings.
An independently specified geometric coefficient recipient
must still receive a map from $\mathcal P_\Gamma$ preserving
the required excess representation, tangent connection and
vacuum density. For noninvertible root substitutions, it
must give the transfer or supported correspondence and its
action on the changing excess. Compatibility is required with
successive substitutions, regular/localized/supported triangles
and field attachments.

The variable equation map and
\eqref{v4-arith:curve:eq:cluster-pullback} are concrete tests.
The ordinary target cannot satisfy the excess requirement.
The constructed derived coefficient target meets it for the
constant equation and invertible coordinate operations above.

\section{The analytic Green difference equation}
\label{v4-arith:curve:sec:analytic}

Let $q>1$ and let $Z(t)$ be rational with $Z(0)=1$. The product
\[
 K(t)=\prod_{j\geq1}Z(q^{-j}t)
\]
defines a meromorphic function on $\mathbb C$, regular at zero
with $K(0)=1$. On a compact set, its tail factors are
$1+O(q^{-j})$ uniformly. Sufficiently far along the tail
they have neither zeros nor poles. The product and its
derivatives therefore converge normally there. The finite
initial product is rational.

Reindexing gives $K(qt)=Z(t)K(t)$.
The normalized holomorphic solution near zero is unique:
the quotient of two such solutions has $H(qt)=H(t)$,
and each positive-degree Taylor coefficient is killed by
$q^n-1$. For torsion Hall data this is the analytic function
with formal expansion $K_G$.

Let $D_Z\subset\mathbb C^\times$ be the finite set of zeros
and poles of $Z$, and put
$U=\mathbb C^\times\setminus q^{\mathbb Z}D_Z$.
Define
\[
 Z_n(t)=
 \begin{cases}
  \prod_{j=0}^{n-1}Z(q^jt),&n>0,\\
  1,&n=0,\\
  \prod_{j=n}^{-1}Z(q^jt)^{-1},&n<0.
 \end{cases}
\]
Then $Z_{m+n}(t)=Z_m(q^nt)Z_n(t)$.
These factors also define the rank-one rational difference
module with relation $\sigma e=Z(t)e$, where
$\sigma f(t)=f(qt)\sigma$. Evaluation $e\mapsto K$ is a
solution in meromorphic functions.

\begin{proposition}\label{v4-arith:curve:prop:difference-line}
The action $(t,v)\mapsto(q^nt,Z_n(t)v)$ descends the trivial
line on $U$ to a holomorphic line on $U/q^{\mathbb Z}$.
The connection $d-d\log K$ descends with it, and $K$ is a
nowhere-zero horizontal section. For
$Z(t)=((1-t)(1-qt))^{-1}$ this analytic trivialization cannot
be obtained from a rational function.
\end{proposition}
\begin{proof}
The action of $q^{\mathbb Z}$ on $\mathbb C^\times$ is free
and properly discontinuous: logarithmic radius changes by
$\log q$. Its zero/pole orbits are locally finite.
Every $Z_n$ is a holomorphic unit on $U$, so the cocycle gives
transition isomorphisms. In logarithmic coordinates the quotient
is $\mathbb C/(2\pi i\mathbb Z+\log q\,\mathbb Z)$ with finitely
many points removed.

The difference equation gives
$q^*(d\log K)=d\log K+d\log Z$, the connection transformation.
The product has no zeros or poles on $U$, and the same
equation makes $K$ the stated section.

If a rational $R$ satisfied $R(qt)/R(t)=Z(t)$, write
$r_b=\operatorname{ord}_b R$. The divisor of its ratio has
order $r_{qb}-r_b$ at $b$. Its sum on each nonzero $q$-orbit
is zero because the divisor of $R$ has finite support.
For the displayed $Z$, the orbit $q^{\mathbb Z}$ instead has
two simple poles, at $1,q^{-1}$, and no zeros. Its sum is
$-2$, a contradiction.
\end{proof}

Thus this line with connection is analytically trivial on the
punctured elliptic quotient. The rational difference equation
has not produced nontrivial monodromy there.
Nor is the quotient by $q^{\mathbb Z}$ a quotient by every
Frobenius ratio: a larger gain group need not act properly
discontinuously.

Between polynomial Hall states, any finite product of the
native Green half-vertices has a matrix element equal to a
finite Laurent polynomial times factors $K(w/z)$.
Normal ordering gives one factor for each annihilation--creation
pair. The remaining matrix element is finite because its two
test states have finite weight. The product construction
continues these expressions meromorphically away from their
locally finite pole divisors. Both bracketings give the same
expression, proving associativity on these matrix elements.
This domain statement asserts no boundedness of an unsmeared
half-vertex on the entire Green Hilbert space.

\section{Hall mass and positive states}
\label{v4-arith:curve:sec:states}

Put $h_m=\sum_{|\lambda|=m}u_\lambda$, $h_0=1$, and
$\mathcal H_Q(t)=\sum_{m\geq0}h_mt^m$.
The group-like Hall series satisfies
\begin{equation}\label{v4-arith:curve:eq:hall-series}
 \mathcal H_Q(t)=\exp\left(\sum_{r\geq1}\frac{p_rt^r}{r}\right).
\end{equation}
Indeed the coefficient of $u_M\otimes u_N$ in
$\Delta\sum_Lu_L$ is
\[
 \sum_L\frac{a_Ma_N}{a_L}F^L_{M,N}
 =\frac{|\operatorname{Ext}^1(M,N)|}{|\operatorname{Hom}(M,N)|}
 =1.
\]

The first equality is orbit--stabilizer for extension classes,
and the last uses the zero torsion Euler form.
Hence $\mathcal H_Q$ is group-like and its logarithm is
primitive. In degree $r$ it is a multiple of $p_r$.

Its indecomposable coefficient is
$\sum_{j=1}^r(-1)^{j-1}\binom{r-1}{j-1}/j=1/r$:
each composition of $r$ gives the unique corresponding flag
in the indecomposable module. The normalization of $p_r$
proves \eqref{v4-arith:curve:eq:hall-series}.

In the nonnegative weight completion set
\[
 v_s=\prod_x\mathcal H_{q^{d_x}}(s^{d_x})
     =\sum_{[M]}s^{\deg M}u_M,
 \qquad \deg M=\sum_xd_x\,\operatorname{length}M_x.
\]
The sum ranges over finite-support torsion objects.
Their Green squared norm is $a_M^{-1}$.

\begin{theorem}\label{v4-arith:curve:thm:mass}
Suppose $N_n\leq C_0q^n$. For $0\leq r<1$,
\[
 K_G(r)=\sum_{[M]}\frac{r^{\deg M}}{a_M}
       =\exp\left(\sum_{n\geq1}\frac{N_nr^n}{n(q^n-1)}\right)
       <\infty .
\]
Thus $v_s$ belongs to the Green Hilbert completion for
$|s|^2=r<1$, with squared norm $K_G(r)$.
For $0<r<1$ the weights
$\mu_r(M)=r^{\deg M}/(a_MK_G(r))$ form a probability measure.
For $0<r<q^{-1}$,
\begin{equation}\label{v4-arith:curve:eq:likelihood}
 \frac{\mu_{qr}(M)}{\mu_r(M)}
   =\frac{q^{\deg M}}{Z_C(r)},\qquad
 \mathbb E_{\mu_r}[q^{\deg M}]=Z_C(r).
\end{equation}
\end{theorem}
\begin{proof}
The factorial primitive norms give
\[
 \|\mathcal H_Q(s)\|_G^2
 =\exp\left(\sum_{r\geq1}\frac{|s|^{2r}}{r(Q^r-1)}\right).
\]
Multiply over the closed points and write $n=rd_x$.
The resulting logarithm is the one in the theorem.
The growth assumption bounds its summands by a constant
times $r^n/n$, so it converges for $r<1$.
Orthogonality in the Hall basis gives the mass sum.
At real nonnegative $r$, partial products and partial sums
increase to their common finite limit.

For $r<q^{-1}$, both series converge. Subtracting their
logarithms gives
\[
 \frac{K_G(qr)}{K_G(r)}=Z_C(r).
\]
Substitution into the two densities proves their likelihood ratio.
Summing proves the expectation.
\end{proof}

A finite map $C\to\mathbb P^1$ gives the required growth
bound by bounding every geometric fiber by its degree.
This argument uses no spectral purity. At $r=0$ the measure
is concentrated on the zero object.

\subsection{Local independence and finite support}

Write
$K_x(r)=\sum_\lambda r^{d_x|\lambda|}/a_\lambda(q^{d_x})$.
The local measures
$r^{d_x|\lambda|}/(a_\lambda K_x(r))$ are independent and
their product gives $\mu_r$. It introduces no infinite-support
objects almost surely, since
\[
 \sum_x\Pr(M_x\ne0)
 =\sum_x(1-K_x(r)^{-1})
 \leq\sum_x\log K_x(r)=\log K_G(r)<\infty .
\]
The first Borel--Cantelli lemma leaves only finitely many
nonzero local factors almost surely. Independence is present
in the construction, though this last implication does not
require it.

The unit vector $\Psi_s=K_G(|s|^2)^{-1/2}v_s$ defines
\[
 \omega_s(P)=\langle\Psi_s,P\Psi_s\rangle_G
\]
on finite polynomials in finitely many creators and Green
annihilators. Every selected mode has a finite coherent
parameter, so every such moment exists. The mean is
\[
 D_{x,m}v_s=\frac{s^{md_x}}{q^{md_x}-1}v_s.
\]
Normal ordering gives all correlations from this mean and
the unchanged Green commutators.
In particular $\omega_s(P^{*G}P)\geq0$.
Infinite smearings require their own test functions and
convergence domain.

\subsection{The degree Hamiltonian and the ordinary trace}

Let $L_0u_M=(\deg M)u_M$ with its closed self-adjoint domain,
specified by square summability after multiplication by
$\deg M$. Its nonnegative integer eigenvalues have the finite
Hall weight multiplicities. Equation~\eqref{v4-arith:curve:eq:likelihood}
becomes the quadratic-form expectation
\[
 \|q^{L_0/2}\Psi_s\|_G^2=Z_C(|s|^2),\qquad q|s|^2<1.
\]
This requires the form domain of $q^{L_0}$. Membership in
the smaller operator domain has not been assumed.
The degree spectrum has been specified independently of zeta zeros.

The ordinary Gibbs trace instead counts orthonormal states once:
\begin{equation}\label{v4-arith:curve:eq:gibbs}
 \operatorname{Tr}_{\mathcal F_C}(e^{-\beta L_0})
 =\prod_{x,r\geq1}(1-e^{-\beta rd_x})^{-1}
 =\prod_{r\geq1}Z_C(e^{-\beta r}),\qquad\beta>\log q .
\end{equation}
Its logarithm is
$\sum_{n\geq1}N_n/[n(e^{\beta n}-1)]$, which converges under
the growth assumption. The Gibbs density is positive and
trace class. In contrast, $K_G(r)=\prod_{j\geq1}Z_C(q^{-j}r)$
is the automorphism-weighted mass and the norm of a coherent
vector. The Green norm of the unnormalized basis does not
insert automorphism weights into the ordinary trace.
Neither positive functional has been identified with the
complete signed Weil functional.

\subsection{A stochastic extension product}

For fixed torsion objects $M,N$, choose an extension class
uniformly in $\operatorname{Ext}^1(M,N)$. Its middle term has law
\[
 P(L\mid M,N)
 =\frac{|\operatorname{Ext}^1(M,N)_L|}
        {|\operatorname{Ext}^1(M,N)|}
 =\frac{a_Ma_N}{a_L}F^L_{M,N}.
\]

The weights sum to one because the torsion Euler form is zero.
In the rescaled basis $v_M=a_Mu_M$, Hall multiplication reads
$v_Mv_N=\sum_LP(L\mid M,N)v_L$. Consequently
\[
 \sum_LP(L\mid M,N)P(K\mid L,P)
 =\sum_LP(L\mid N,P)P(K\mid M,L).
\]
Both sides are $a_Ma_Na_P/a_K$ times the same three-step
flag count. Higher regroupings are the same higher flag
counts, so this gives all finite stochastic associativity
identities. An enhanced primitive object still requires
its cochain or sheaf-level comparison.

\section{Prime-indexed Macdonald and Hall constructions}
\label{v4-arith:etingof}

The decomposition by closed points suggests a comparison with
symmetric functions at each prime. This construction has two
independent choices: the set of tensor factors and the parameters
inside each factor. A single pair of parameters cannot equal
several distinct pairs simultaneously. This elementary fact does
not prohibit a family of Macdonald spaces or a map of their
underlying algebras. The relevant obstruction appears only after
specifying which gradings and operations the map must preserve.

\subsection{Pairings and the degree obstruction}
\label{v4-arith:et:macdonald-qt}

Let $K=\mathbb Q(q_p,t_p:p\text{ prime})$ be the field of
rational functions in algebraically independent variables $q_p,t_p$.
Each rational function involves finitely many variables. Put
$\Lambda_{(p)}=K[p_1^{(p)},p_2^{(p)},\ldots]$, with
$\deg p_r^{(p)}=r$.
For a partition $\lambda$, let $m_r(\lambda)$ be the number of
parts equal to $r$, and set
\[
 p_\lambda^{(p)}=\prod_r(p_r^{(p)})^{m_r(\lambda)},\qquad
 z_\lambda=\prod_r r^{m_r(\lambda)}m_r(\lambda)!.
\]

\begin{definition}[Macdonald pairing]
\label{v4-arith:thm:macdonald-norm}
The symmetric bilinear form on $\Lambda_{(p)}$ is specified by
\[
 \langle p_\lambda^{(p)},p_\mu^{(p)}\rangle_p
 =\delta_{\lambda\mu}z_\lambda
   \prod_{j=1}^{\ell(\lambda)}
        \frac{1-q_p^{\lambda_j}}{1-t_p^{\lambda_j}}.
\]
This is the normalization of Macdonald,
\emph{Symmetric Functions and Hall Polynomials}, VI, equation (1.5).
\end{definition}

\begin{definition}
\label{v4-arith:def:per-prime-macdonald}
Define
\[
 \Lambda^{\mathrm{prim}}=\bigotimes_p'\Lambda_{(p)}
   =K[p_r^{(p)}:p\text{ prime},\ r\geq1].
\]
The prime denotes the algebraic tensor product with unit $1$ in
all but finitely many factors. Give it the tensor product pairing
and total degree $\deg p_r^{(p)}=r$.
\end{definition}

The form is nondegenerate. Every polynomial lies in a finite tensor
product, where its power-sum monomials have nonzero diagonal norms.
For real parameters $0<q_p,t_p<1$, the same diagonal formula
defines a positive Hermitian form on the complex polynomial space,
with conjugation in the first argument. Positivity holds for every finite linear combination. It
requires no infinite tensor product completion.

\begin{proposition}\label{v4-arith:obs:domain-mismatch-macdonald}
Let $\Lambda=K[p_1,p_2,\ldots]$, graded by $\deg p_r=r$.
A tensor product of two or more $\Lambda_{(p)}$ admits no
injective degree-preserving $K$-linear map to $\Lambda$.
After forgetting degree, $\Lambda^{\mathrm{prim}}$ and $\Lambda$
are isomorphic as $K$-algebras.
\end{proposition}
\begin{proof}
The degree-one space of $\Lambda$ has dimension one. Each retained
prime contributes an independent degree-one generator to the
source, proving the obstruction. For the second assertion, choose
a bijection $b$ from the pairs $(p,r)$ to the positive integers.
The assignments $p_r^{(p)}\mapsto p_{b(p,r)}$ and their inverse
extend uniquely to mutually inverse polynomial algebra maps.
They need not preserve degree, the pairing or any marked operator.
\end{proof}

Thus $q_p=p^{-s}$ and $t_p=p^{-s-1}$, for real $s>0$, give a
well-defined family of positive pairings. They do not specify one
constant pair $(q,t)$ independent of $p$. An arithmetic comparison
must say whether it preserves these pairings and the prime labels.
Parameter variation alone is not a nonexistence theorem.

\subsection{Adjoints, weights and Macdonald bases}
\label{v4-arith:et:logp-not-in-qt}

Write $c_{p,r}=(1-q_p^r)/(1-t_p^r)$. Let $M_{p,r}$ denote
multiplication by $p_r^{(p)}$. Its adjoint is
\[
 D_{p,r}=r c_{p,r}\partial_{p_r^{(p)}},\qquad
 [D_{p,r},M_{\ell,s}]
 =\delta_{p\ell}\delta_{rs}r c_{p,r}I.
\]
Indeed, the squared norm of $(p_r^{(p)})^m$ is
$m!(r c_{p,r})^m$. Comparing consecutive norms proves adjunction.
This also identifies exactly which coefficient changes when the
pairing changes.

\begin{proposition}\label{v4-arith:thm:free-fermion-match}
At $t_p=q_p$, one has $c_{p,r}=1$. Along $t_p=q_p^2$,
\[
 c_{p,r}=\frac1{1+q_p^r}\longrightarrow1
 \quad\text{as }q_p\longrightarrow0.
\]
Both give the ordinary Heisenberg squared norm $r$ for $p_r^{(p)}$
in the stated specialization or limit. For any positive numbers $w_p$, replacing
$c_{p,r}$ by $w_p^r$ defines another positive form, with
\[
 \|p_\lambda^{(p)}\|^2=z_\lambda w_p^{|\lambda|}.
\]
\end{proposition}
\begin{proof}
The two formulas for $c_{p,r}$ follow by cancellation. Every
power-sum monomial then has the claimed factorial norm.
All its factors are positive in the weighted form.
\end{proof}

The choice $w_p=\log p$ yields
$\|p_r^{(p)}\|^2=r(\log p)^r$.
It is extra metric data. Over the real numbers, the map
$p_r^{(p)}\mapsto w_p^{r/2}p_r^{(p)}$ is an algebra isometry from
this weighted space to the ordinary Hall inner-product space.
It need not commute with the usual translation
$T p_r^{(p)}=r p_{r+1}^{(p)}$: its two orders of application give
factors $w_p^{r/2}$ and $w_p^{(r+1)/2}$.
A conformal comparison must also transport translation and the
chosen conformal vector.

\begin{example}\label{v4-arith:comp:norm-vs-kesten-mckay}
In the ordinary Heisenberg pairing, $p_r^{(p)}/\sqrt r$ has norm
one. This is a statement about a Fock vector. For adjacency on the
$(p+1)$-regular tree, the root spectral measure instead has moments
\[
 \int x^2\,d\nu_p=p+1,\qquad
 \int x^4\,d\nu_p=(p+1)(2p+1).
\]

There are $p+1$ two-step returns. A four-step return either returns
at step two, giving $(p+1)^2$ paths, or first reaches distance two,
giving $(p+1)p$ paths. Odd moments vanish because the tree is
bipartite. These state moments differ from the norms of individual
power-sum vectors. Their Hecke recurrence is the one in
Proposition~\ref{v4-arith:hecke:prop:distance}.
\end{example}

\subsection{The tensor product basis}
\label{v4-arith:et:arithmetic-macdonald}

Let $m_\lambda$ be the monomial symmetric function. In dominance
order, $\mu\leq\lambda$ means
$\sum_{i\leq j}\mu_i\leq\sum_{i\leq j}\lambda_i$ for every $j$,
with equal total size. Macdonald's theorem, VI (4.7), gives the
unique orthogonal basis $P_\lambda(x;q,t)$ over $\mathbb Q(q,t)$
which is monic triangular in $m_\lambda$ for this order.
At specialized parameters the coefficients must remain defined.

\begin{definition}\label{v4-arith:def:arithmetic-macdonald-poly}
For a family of partitions with finite support, put
\[
 P^{\mathrm{arith}}_{(\lambda_p)}
 =\bigotimes_p'P_{\lambda_p}(x^{(p)};q_p,t_p),
 \qquad P_\varnothing=1.
\]
\end{definition}

These tensors form an orthogonal $K$-basis. Both expansion and
linear independence reduce to a finite set of prime factors.
For two primes, the degree-two basis consists of
\[
 \begin{gathered}
 P_{(2)}\otimes1,\quad P_{(1,1)}\otimes1,\quad
 P_{(1)}\otimes P_{(1)},\\
 1\otimes P_{(2)},\quad1\otimes P_{(1,1)}.
 \end{gathered}
\]
A single symmetric-function space has degree-two dimension two,
with basis $p_2,p_1^2$. Symmetric polynomials in two variables also
have dimension two in that degree. The three ordinary monomials
$x_1^2,x_1x_2,x_2^2$ belong to the full polynomial space, whose
symmetric subspace has basis $x_1^2+x_2^2,x_1x_2$.

\subsection{DAHA and its geometric comparisons}
\label{v4-arith:et:no-prime-daha}

The double affine Hecke algebra, abbreviated DAHA, acts on a fixed
number of Laurent variables. This number is independent of the
set of primes indexing a tensor product. The distinction can be
seen directly in its polynomial representation.

\begin{theorem}[Polynomial representation]
\label{v4-arith:thm:DAHA-rank}
Over $K_0=\mathbb Q(q^{1/2},t^{1/2})$, the DAHA of type
$\operatorname{GL}_n$ has a faithful representation on
$K_0[x_1^{\pm1},\ldots,x_n^{\pm1}]$. In the convention of
Schiffmann--Vasserot, arXiv:0905.2555, section 1.4, its operators are
\[
 \begin{split}
 X_i&=x_i,\\
 T_i&=t^{1/2}s_i+
       \frac{t^{-1/2}-t^{1/2}}{x_i/x_{i+1}-1}(s_i-1),\\
 Y_i&=T_i\cdots T_{n-1}\omega T_1^{-1}\cdots T_{i-1}^{-1}.
 \end{split}
\]
Here $s_i$ exchanges $x_i,x_{i+1}$,
$\partial_i f(x)=f(x_1,\ldots,qx_i,\ldots,x_n)$ and
$\omega=s_{n-1}\cdots s_1\partial_1$.
The spherical subalgebra is $e\mathbb H_ne$, where $e$ is the
normalized finite Hecke symmetrizer.
\end{theorem}

This is the standard faithful polynomial representation theorem.
The displayed divided difference preserves Laurent polynomials,
since $(s_i-1)f$ is divisible by $x_i-x_{i+1}$.
On symmetric polynomials the first spherical operator becomes
\[
 \sum_{i=1}^n
  \left(\prod_{j\ne i}
       \frac{t^{-1/2}x_i-t^{1/2}x_j}{x_i-x_j}\right)\partial_i.
\]
Its eigenvectors are $P_\lambda(x;q,t^{-1})$, in this convention,
with eigenvalues $t^{-(n-1)/2}\sum_iq^{\lambda_i}t^{i-1}$.
The inverse in the second Macdonald parameter is part of the
normalization.

\begin{remark}\label{v4-arith:obs:DAHA-Adams-incompatible}
The Weyl group acts inside each alphabet. Adams operations can
also be defined factor by factor by
$\psi_m(p_r^{(p)})=p_{mr}^{(p)}$, fixing the coefficient field.
Operators on distinct tensor factors commute. These definitions
pose no incompatibility between prime labels and DAHA rank.
An intertwiner between the two operator families remains additional
data, as does a geometric interpretation of the factorwise
Adams operations.
\end{remark}

Two geometric results specify which Hall algebra enters this
comparison. Let $T=(\mathbb C^\times)^2$ act on the affine plane,
and let $L=\bigoplus_{n\geq0}K_T(\operatorname{Hilb}^n\mathbb A^2)$
with equivariant parameters localized to their fraction field.
Let $H$ be the convolution algebra generated by the tautological
line-bundle powers on adjacent nested Hilbert schemes and the
Adams operations of the tautological bundles.

\begin{theorem}[Schiffmann--Vasserot]
\label{v4-arith:thm:SV-at-q-to-1}
In arXiv:0905.2555, Theorem 3.1 identifies $H$ with the central
specialization $\mathcal E_c$ of the elliptic Hall algebra at
$c=(1,q^{1/2}t^{1/2})$. Its representation on $L$ is the standard
symmetric-function representation.

The stable centrally extended degenerate spherical DAHA
$\mathrm{SH}^c$ of arXiv:1202.2756 has an integral form at
$\kappa=1$. For its specialization $\mathrm{SH}^c_1$,
Theorem F.10 gives
\[
 \mathrm{SH}^c_1\cong
 U(W_{1+\infty})\otimes_{\mathbb C}
                  \mathbb C[c_1,c_2,\ldots].
\]
Here $W_{1+\infty}$ is the universal central extension of the Lie
algebra of regular differential operators on the circle.
\end{theorem}

The second construction uses the degeneration
$q,t\longrightarrow1$ with $t=q^{-\kappa}$, together with its
integral form and central parameters. Its result is an enveloping
algebra. A bare substitution $q=1$ in a rational shuffle formula
does not specify this degeneration. Neither geometric theorem
constructs the arithmetic Hochschild trace class
$\mathcal V^{\mathrm{prim}}$.

\subsection{A rational shuffle construction}
\label{v4-arith:et:shuffle-uncoupled}

One can construct a prime-indexed shuffle algebra without making
the geometric identification. In the field $K$ above, put
\[
 \omega_p(x,y)=
 \frac{(1-q_px/y)(1-t_p^{-1}x/y)}
      {(1-x/y)(1-q_pt_p^{-1}x/y)}.
\]
For $n\geq0$, let $\mathscr S_{p,n}=K(x_1,\ldots,x_n)^{S_n}$,
with $\mathscr S_{p,0}=K$. For symmetric rational functions
$f\in\mathscr S_{p,r}$ and $g\in\mathscr S_{p,s}$, define
\[
 (f*g)(x_1,\ldots,x_{r+s})=
 \sum_{\substack{I\sqcup J=\{1,\ldots,r+s\}\\|I|=r}}
 f(x_I)g(x_J)\prod_{i\in I,\,j\in J}\omega_p(x_i,x_j).
\]
This is the sum over subsets, with no factorial normalization.

\begin{definition}\label{v4-arith:def:arithmetic-hall}
Set $\mathscr S_p=\bigoplus_{n\geq0}\mathscr S_{p,n}$ and
$\mathscr S^{\mathrm{prim}}=\bigotimes_p'\mathscr S_p$, with
unit $1$ in degree zero. Separately, a specified family of
spherical elliptic Hall algebras $\mathrm{SH}(q_p,t_p)$ defines
\[
 \mathbb H^{\mathrm{arith}}=
       \bigotimes_p'\mathrm{SH}(q_p,t_p).
\]
All factors are extended to the common coefficient field.
\end{definition}

\begin{proposition}\label{v4-arith:thm:shuffle-uncoupled}
The displayed shuffle product is associative and unital.
In $\mathscr S^{\mathrm{prim}}$, factors at distinct primes commute,
and their cross-prime shuffle coefficient is one. The same
commutation holds in $\mathbb H^{\mathrm{arith}}$ by its tensor
product definition.
\end{proposition}
\begin{proof}
Both orders of a triple shuffle sum over ordered decompositions
$I\sqcup J\sqcup L$. Their summand is $f(x_I)g(x_J)h(x_L)$
times the kernels for $I\times J$, $I\times L$ and $J\times L$.
Each pair occurs once in both expressions. The scalar kernels
commute, proving associativity. The empty block proves the unit
law. Multiplication in distinct tensor factors introduces no kernel.
\end{proof}

This proves a property of the specified tensor product. It does
not forbid other arithmetic constructions with cross-prime
operations. Identifying a shuffle subalgebra with a positive
elliptic Hall half requires its presentation and wheel conditions.
Recovering the full algebra also requires the negative and Cartan
parts and their relations. Associativity of the rational kernel
alone proves none of these identifications.

\begin{proposition}\label{v4-arith:thm:parameter-non-uniformity}
The algebraic prime tensor space carries the nondegenerate
Macdonald pairing, the tensor Macdonald basis and the factorwise
Adams operations above. Its independent parameters are compatible
with a prime-indexed family of shuffle or spherical Hall algebras.
For the rational shuffle construction, there are no mixed-prime
coefficients beyond tensor multiplication.
\end{proposition}
\begin{proof}
Each tensor and each operation uses finitely many factors.
Nondegeneracy and the basis assertion reduce to those factors.
The Adams substitutions preserve polynomial multiplication and
satisfy $\psi_m\psi_n=\psi_{mn}$. The final assertion is
Proposition~\ref{v4-arith:thm:shuffle-uncoupled}.
\end{proof}

An arithmetic realization must still construct a map to the
Hochschild trace carrier and state which operators, pairing,
grading and trace functional it preserves. The polynomial space
and its prime labels supply a possible domain for such a map.
They do not identify its target.


\section{A geometric Hall kernel on finite differential jets}
\label{v4-arith:curve:sec:jets}

The curve construction uses a finite field cardinality $q$.
The following example uses a separate nilpotent parameter
$\epsilon$. It compares the equivariant cohomological Hall
algebra of the three-loop quiver with \emph{zero potential}
to an explicitly constructed field algebra on differential jets.
The critical algebra for $\operatorname{Tr}(Z[X,Y])$ has
different coefficients and requires a further oriented map.

Let $k=\mathbb Q(a,b)$, $h=a+b$, and
$R_N=k[\epsilon]/(\epsilon^N)$, $N\geq1$.
Give the three arrows weights $\epsilon a,\epsilon b,-\epsilon h$.
This is a base change of the two-dimensional equivariant
parameter ring. It inverts neither $\epsilon$ nor $h\epsilon$.
Put
\begin{equation}\label{v4-arith:curve:eq:euler-kernels}
 \begin{aligned}
 \ell&=\epsilon h,& P(s)&=s^2-\ell^2,\\
 K_0(s)&=\frac{(s+\epsilon a)(s+\epsilon b)(s-\ell)}s,
 &E_0(s)&=K_0(s)^{-1},\\
 E_{\mathrm{pl}}(s)&=\frac{s(s+\ell)}{(s+\epsilon a)(s+\epsilon b)},
 &r(s)&=\frac{E_{\mathrm{pl}}(s)}{E_{\mathrm{pl}}(-s)}.
 \end{aligned}
\end{equation}

Then $E_0=E_{\mathrm{pl}}/P$ and $E_0(s)/E_0(-s)=r(s)$.
Expand in $\epsilon$ first, with the indicated spectral
variable inverted.

The charge-$d$ coefficient algebra is
$A_d=R_N[x_1,\ldots,x_d]^{\mathfrak S_d}$.
The representation space before base change is
$\operatorname{End}(\mathbb C^d)^3$. Its equivariant cohomology
has this polynomial description.

The invariant-subspace correspondence gives
\begin{equation}\label{v4-arith:curve:eq:geometric-product}
 (F*G)(X)=
 \sum_{\substack{I\sqcup J=[d+e]\\|I|=d}}
 F(X_I)G(X_J)\prod_{i\in I,j\in J}K_0(x_j-x_i).
\end{equation}
There is no factorial divisor in this proper flag pushforward.

At a coordinate flag, the three moving arrow Euler factors
are $y-x+\epsilon a,y-x+\epsilon b,y-x-\ell$.
The flag tangent factor is $y-x$. Their ratio is $K_0$.
This is the shuffle calculation of
Theorem~\ref{v4-arith:hall:thm:coha-shuffle} with these weights.

\begin{proposition}\label{v4-arith:curve:prop:geometric-hall}
The product \eqref{v4-arith:curve:eq:geometric-product} is
polynomial and associative. Its inverse-Euler spectral
coefficient map is
\begin{equation}\label{v4-arith:curve:eq:hall-split}
 \Delta^0_{d,e}(s)F
 =\prod_{i,j}E_0(s+y_j-x_i)\,F(x-s,y)
 \in(A_d\otimes_{R_N}A_e)((s^{-1})).
\end{equation}
It is rationally coassociative. In charge split $(1,1)$,
\[
 1*1=2\bigl((y-x)^2+\epsilon^2(ab-h^2)\bigr),\qquad
 \Delta^0_{1,1}(s)1=E_0(s+y-x).
\]
\end{proposition}
\begin{proof}
Proper pushforward is defined before equivariant localization,
so the shuffle sum is polynomial. Directly, residues on each
possible root diagonal cancel between the two flags exchanged
there. Either bracketing of three factors sums over the same
ordered three-block partitions with one $K_0$ per ordered
cross-block pair. This proves associativity with the displayed
normalization.

The inverse Euler factors and translation of the first block
give \eqref{v4-arith:curve:eq:hall-split}.
For a triple split, both composites have one factor for each
pair, at parameters $s,t,s+t$, and evaluate the input at
$(x-s-t,y-t,z)$. They agree before Laurent expansion.
Finally $K_0(v)+K_0(-v)$ gives the displayed product in
charge two.
\end{proof}

With roots and $\epsilon$ of degree two, a root polynomial
of degree $m$ in charge $d$ has shifted Hall degree
$2m-2d^2$. The product is homogeneous and all charges have
even parity. This grading is distinct from the cochain
grading, which places the differential Hopf states below
in degree zero.

There are two reasons to change the polynomial carrier before
constructing these fields. Before nilpotent base change, take
the three arrow weights to be $1,1,-2$. Then
\[
 r(s)=\frac{(s+2)(s-1)^2}{(s-2)(s+1)^2},
 \qquad\operatorname{Res}_{s=2}r(s)\,ds=\frac49 .
\]

Its expansion has an infinite nonzero tail.
A coefficientwise injective marking preserves that tail.
Thus the unchanged exchange cannot become a finite Laurent
polynomial on the same pair of states. Completing only the charge grading leaves this
charge-two obstruction unchanged, and a modification supported
only at $s=0$ cannot cancel the pole at $s=2$.

Also, if $T^Ma=0$ and
$Y(Ta,z)b=\partial_zY(a,z)b$, then
$\partial_z^MY(a,z)b=0$.
Over a $\mathbb Q$-algebra the falling factorial of a
negative integer is a unit.
A lower-truncated Laurent series with this property is
therefore a polynomial of degree less than $M$.
The root derivative $-\sum_i\partial_{x_i}$ is locally
nilpotent on $A_d$. Finite jets reverse this behavior:
the transpose derivative raises jet order.

\subsection{A difference primitive and its Hopf form}

We first realize the normalized factor $E_{\mathrm{pl}}$.
There is a unique formal series
$F(z)\in(\epsilon/z)k[[\epsilon/z]]$ satisfying
\begin{equation}\label{v4-arith:curve:eq:difference-primitive}
 F(z-\ell)-F(z)=\log E_{\mathrm{pl}}(-z).
\end{equation}

To prove it, write $t=\epsilon/z$, $F=\sum_{n\geq1}f_nt^n$,
and
$g_j=(a^j+b^j-h^j)/j$.
The right side is $\sum_{j\geq2}g_jt^j$. Comparison of
coefficients gives the recurrence
\[
 f_{j-1}=
 \frac{g_j-\sum_{n=1}^{j-2}
                f_n\binom{j-1}{j-n}h^{j-n}}
      {(j-1)h}\qquad(j\geq2).
\]

Its denominator is a unit in $k$, and no division by
$\epsilon$ occurs. This construction uses $h^{-1}$ and
has not been specialized to $h=0$.
The identity
$E_{\mathrm{pl}}(z-\ell)=E_{\mathrm{pl}}(-z)$ shows that
$-F(-z)$ solves the same equation. Uniqueness makes $F$ odd.

With $\varpi(z)=\exp F(z)$, this gives
\begin{equation}\label{v4-arith:curve:eq:varpi}
 \varpi(z)\varpi(-z)=1,\qquad
 \frac{\varpi(z-\ell)}{\varpi(z)}=E_{\mathrm{pl}}(-z),\qquad
 \frac{\varpi(z)^2}{\varpi(z-\ell)\varpi(z+\ell)}=r(z).
\end{equation}

For example,
$F(z)=-(ab/h)\epsilon z^{-1}
 +(a^2b^2/(6h))\epsilon^3z^{-3}+O(\epsilon^5)$.
Reduce all these series modulo $\epsilon^N$.

Take the commutative, cocommutative Hopf algebra
\[
 H_N=R_N[g,g^{-1},b_0,b_1,\ldots],
\]
where $g$ is group-like, the $b_i$ are primitive,
$\varepsilon(g)=1$, $\varepsilon(b_i)=0$,
$Sg=g^{-1}$ and $Sb_i=-b_i$.
The derivation $T_Hg=gb_0$, $T_Hb_i=b_{i+1}$ is also a
coderivation and satisfies $\varepsilon T_H=0$.
Every algebraic coproduct is finite.

For integers $m,n$ and finite lists of formal variables $u_i,v_j$,
define a bilinear form by differentiation at zero in those
variables:
\begin{align*}
 \rho^H_z\left(g^me^{\sum_i u_ib_i},g^ne^{\sum_jv_jb_j}\right)
 &=\exp B^H_z,\\
 B^H_z&=mnF+n\sum_i u_iF^{(i+1)}
       +m\sum_j(-1)^{j+1}v_jF^{(j+1)}\\
 &\quad+\sum_{i,j}(-1)^{j+1}u_iv_jF^{(i+j+2)} .
\end{align*}
Here $F^{(i)}=\partial_z^iF$.
Differentiation, rather than an unnormalized coefficient,
specifies the factorial convention on monomials.
Every coefficient of $B^H_z$ is divisible by $\epsilon$,
so the form has values in $R_N[z,z^{-1}]$, including on
negative powers of $g$.

The identities needed below can be stated explicitly.
For a Hopf algebra with translation $T$, write
$a(z)=e^{zT}a$, and sum over coproduct subscripts. They are
\begin{equation}\label{v4-arith:curve:eq:six-identities}
 \begin{aligned}
 \rho_z(1,a)&=\varepsilon(a)=\rho_z(a,1),\\
 \rho_z(ab,c)&=\rho_z(a,c_1)\rho_z(b,c_2),\\
 \rho_z(a,bc)&=\rho_z(a_1,c)\rho_z(a_2,b),\\
 \rho_z(a_1,b_1)a_2(z)b_2&=b_1a_1(z)\rho_z(a_2,b_2),\\
 \rho_z(Ta,b)&=\partial_z\rho_z(a,b)=-\rho_z(a,Tb),\\
 \rho_z(a_1,b_1)\rho_{-z}(b_2,a_2)&=\varepsilon(a)\varepsilon(b).
 \end{aligned}
\end{equation}

The form $\rho^H$ satisfies all six.
The formal exponential inputs are group-like, and $B^H_z$
is additive in each input list. This proves the two
multiplicativities after differentiation.
A zero list gives the unit rule.

Commutativity and cocommutativity give the fourth identity.
The first-input translation rule on $(g,g)$ is
\[
 \rho^H_z(T_Hg,g)=\varpi F'=\varpi'.
\]
The second-input derivative has the opposite sign. On the $b_i$, it follows
by differentiating the displayed entries. Derivation and
coderivation extend it to products and inverses.

Finally $F^{(i)}(-z)=(-1)^{i+1}F^{(i)}(z)$ makes the two
exponents cancel when inputs are exchanged and $z$ is
replaced by $-z$. This proves the last, convolution identity.

\subsection{Restoring the even Euler factor}

The factor $P^{-1}$ in $E_0$ changes the contraction while
canceling from the exchange. To retain it, use a second
commutative, cocommutative Hopf algebra
\[
 B_N=R_N[U,U^{-1},V,V^{-1},c_{\alpha,i}:
                        \alpha=1,2,\ i\geq0].
\]
Let $g_1=U,g_2=V$ be group-like and the $c_{\alpha,i}$ primitive.
Set $T_Bg_\alpha=g_\alpha c_{\alpha,0}$,
$T_Bc_{\alpha,i}=c_{\alpha,i+1}$, with the corresponding
group-like and primitive antipodes and counits.
Write
$J=\left(\begin{smallmatrix}0&-1\\1&0\end{smallmatrix}\right)$
and $L(z)=P'(z)/P(z)$.

The form is specified by
\[
 \begin{aligned}
 \rho^B_z(g_\alpha,g_\beta)&=P(z)^{J_{\alpha\beta}},\\
 \rho^B_z(c_{\alpha,i},g_\beta)&=J_{\alpha\beta}L^{(i)}(z),\\
 \rho^B_z(g_\alpha,c_{\beta,j})
    &=(-1)^{j+1}J_{\alpha\beta}L^{(j)}(z),\\
 \rho^B_z(c_{\alpha,i},c_{\beta,j})
    &=(-1)^{j+1}J_{\alpha\beta}L^{(i+j+1)}(z).
 \end{aligned}
\]
Precisely, its value on
$g^me^{\sum p_{\alpha,i}c_{\alpha,i}}$ and
$g^ne^{\sum q_{\beta,j}c_{\beta,j}}$ is
$P^{m^tJn}\exp B^B_z$.
Here $m,n\in\mathbb Z^2$, $g^m=U^{m_1}V^{m_2}$,
and the auxiliary variable lists are finite. Set
\[
 \begin{split}
 B^B_z=\sum_{\alpha,\beta}J_{\alpha\beta}\biggl(
 &n_\beta\sum_i p_{\alpha,i}L^{(i)}
 +m_\alpha\sum_j(-1)^{j+1}q_{\beta,j}L^{(j)}\\
 &+\sum_{i,j}(-1)^{j+1}p_{\alpha,i}q_{\beta,j}L^{(i+j+1)}
 \biggr).
 \end{split}
\]

Again differentiate in the auxiliary variables at zero.
For each coefficient only finitely many exponential terms
contribute. Integral powers of $P$ and derivatives of $L$
are finite Laurent polynomials modulo $\epsilon^N$.

This form also satisfies
\eqref{v4-arith:curve:eq:six-identities}.
Additivity of $B^B$ and bilinearity of $m^tJn$ prove
multiplicativity. Units and commutation follow as before.
Translation starts from $(P^J)'=JP^JL$.

For the last identity, $P$ is even, $J$ is alternating,
and $L^{(i)}(-z)=(-1)^{i+1}L^{(i)}(z)$.
The powers of $P$ and all exponential terms cancel under
reversal. This proof does not require $L$ to be divisible
by $\epsilon$.

The shear
\[
 \sigma U=UV,\quad \sigma V=V,\quad
 \sigma c_{1,i}=c_{1,i}+c_{2,i},\quad
 \sigma c_{2,i}=c_{2,i}
\]
is a Hopf automorphism commuting with $T_B$ and preserving
the form. Its matrix
$A=\left(\begin{smallmatrix}1&0\\1&1\end{smallmatrix}\right)$
satisfies $A^tJA=J$.
Its inverse sends $U$ to $UV^{-1}$ and $c_{1,i}$ to
$c_{1,i}-c_{2,i}$, fixing $V$ and the $c_{2,i}$.

On
$\widetilde H_N=H_N\otimes_{R_N}B_N$
take the product form, total translation $T$, and
\[
 \vartheta=e^{\ell T_H}\otimes\sigma .
\]

The exponential is finite because $\ell^N=0$.
It is a Hopf automorphism preserving the form:
simultaneous translation differentiates a scalar value by
$\partial_z-\partial_z=0$.
Such an automorphism is the \emph{level} used in this
algebraic construction. It is additional data on the
specified Hopf carrier.

Put $G=gU$, $G(v)=e^{vT}G$ and $J_m=T^mG$.
The subalgebra
$W_N^+=R_N[J_0,J_1,\ldots]$ is freely polynomial.
After $U=V=1,c_{\alpha,i}=0$, the $J_m$ become $T_H^mg$.
Their ratios by $g$ are triangular polynomials in
$b_0,\ldots,b_{m-1}$ with coefficient one on $b_{m-1}$.
Localization at $g$ proves algebraic independence.

Give each $J_m$ charge one and denote charge $d$ by
$W_{N,d}^+$.
The ambient charge assigns $g$ degree one and
$U,V,b_i,c_{\alpha,i}$ degree zero.
It restricts to this grading and is preserved by $\vartheta$.

\subsection{The field and its charge bounds}

For any level $\theta$ preserving the form, define
\[
 \delta_\theta(a)=a_2\otimes a_3S^{-1}(\theta a_1)
                =a_{(0)}\otimes a_{(1)}.
\]
The field and exchange on the algebraic Hopf carrier are
\begin{equation}\label{v4-arith:curve:eq:hopf-field}
 \begin{aligned}
 Y_\theta(a,z)b&=a_1(z)b_{(0)}
                    \rho_{-z}(b_{(1)},a_2),\\
 \mathcal S_\theta(z)(a\otimes b)
   &=a_{(0)}\otimes b_{(0)}\rho_z(a_{(1)},b_{(1)}).
 \end{aligned}
\end{equation}
All sums are finite.

The unit rule gives
$Y_\theta(1,z)=I$ and $Y_\theta(a,z)1=e^{zT}a$.
The translation rule and total coderivation rule for
$\delta_\theta$ give covariance.
For this commutative Hopf algebra, we check the remaining
field identities on group-like generating series.
Auxiliary differentiation then gives them on all states.

\begin{theorem}\label{v4-arith:curve:thm:jet-fields}
The fields at level $\vartheta$ preserve $W_N^+$, add
creation charges, and satisfy
\begin{align}
 Y_\vartheta(G,z)G&=E_0(-z)G(z)G,\qquad
 \mathcal S_\vartheta(z)(G\otimes G)=r(z)G\otimes G,
       \label{v4-arith:curve:eq:jet-generator}\\
 Y_\vartheta\left(\prod_iG(u_i),z\right)\prod_jG(v_j)
 &=\prod_iG(z+u_i)\prod_jG(v_j)
       \prod_{i,j}E_0(-z-u_i+v_j).
       \label{v4-arith:curve:eq:jet-blocks}
\end{align}
The exchange of the two blocks multiplies by
$\prod_{i,j}r(z+u_i-v_j)$.
Every finite product has a common rational expansion with
the field associativity and exchange identities.
For charges $d,e$ and total input jet order $M$, a pole
bound is $2de+N-1+M$.
\end{theorem}
\begin{proof}
The coaction of $G$ is $G\otimes C$, where
\[
 C=G\vartheta(G)^{-1}
   =\frac{g}{e^{\ell T_H}g}\,V^{-1}.
\]
Equations~\eqref{v4-arith:curve:eq:varpi} and the auxiliary
form give
$\rho_{-z}(C,G)=E_{\mathrm{pl}}(-z)P(z)^{-1}=E_0(-z)$
and $\rho_z(C,C)=r(z)$.
The auxiliary contribution to the latter is
$\rho^B_z(V^{-1},V^{-1})=1$.
This proves the generator formulas.
Multiplicativity and translation give the block formulas.
Auxiliary derivatives express every coefficient in the
polynomial algebra on the $J_m$.

For three or more blocks, either nesting includes the same
contractions between each pair of elementary generating
states, with the same translated regular product.
This proves weak associativity as equality of expansions
of a common rational germ.

Interchanging adjacent blocks
uses $r(s)E_0(-s)=E_0(s)$.
Reversal uses $r(s)r(-s)=1$.
Three exchanges multiply the same scalar factors, giving
the Yang--Baxter identity, and splitting a block factors
the same product, giving compatibility with a field
insertion.

Differentiating proves these identities on
every polynomial jet input. Clearing the finitely many
pairwise denominators gives the corresponding exchange
locality, with a bound depending on the two inputs.

Each contraction has baseline inverse degree two.
For $de$ pairs the baseline is $2de$. Its further inverse
degree equals its $\epsilon$ degree and is at most $N-1$.
Input differentiation adds at most $M$. This proves the
stated lower Laurent bound and all required truncations.
\end{proof}

Modulo $\epsilon^4$,
\begin{equation}\label{v4-arith:curve:eq:cubic-jet}
 E_0(-z)=z^{-2}+(h^2-ab)\epsilon^2z^{-4}
                   -abh\epsilon^3z^{-5},\qquad
 r(z)=1+2abh\epsilon^3z^{-3}.
\end{equation}
The even factor changes the zeroth-order contraction from
$1$ to $z^{-2}$, while leaving the cubic exchange unchanged.

\subsection{Transposing the Hall coefficient map}

Let $A_d[m]$ denote homogeneous root degree $m$, and put
$A_d^\circ=\bigoplus_m\operatorname{Hom}_{R_N}(A_d[m],R_N)$.
This is the continuous dual of $\prod_mA_d[m]$ with its
degree topology. Pair it with $W_{N,d}^+$ by
\begin{equation}\label{v4-arith:curve:eq:jet-pairing}
 \langle J_{m_1}\cdots J_{m_d},F\rangle
 =\left.\partial_{x_1}^{m_1}\cdots\partial_{x_d}^{m_d}
                           F\right|_{x=0}.
\end{equation}
Monomial symmetric functions and unordered derivative
lists pair diagonally with nonzero factorials. Hence this
pairing is perfect in every degree and gives
$\iota_d:A_d^\circ\cong W_{N,d}^+$.

\begin{theorem}\label{v4-arith:curve:thm:transpose}
For finite-jet inputs of charges $d,e$,
\[
 Y_\vartheta(\iota_d\alpha,z)\iota_e\beta
 =\iota_{d+e}\bigl((\Delta^0_{d,e}(-z))^\circ
                                      (\alpha\otimes\beta)\bigr).
\]
The right side is defined coefficientwise on polynomial
tests and lies in $W_{N,d+e}^+((z))$.
Translation is $T=-D^\circ$ for
$D=-\sum_i\partial_{x_i}$.
The exchange is the transpose of the inverse Hall
multiplier at parameter $-z$.
\end{theorem}
\begin{proof}
Pair \eqref{v4-arith:curve:eq:jet-blocks} with $F$.
The result is
\[
 \prod_{i,j}E_0(-z-u_i+v_j)
     F(z+u_1,\ldots,z+u_d,v_1,\ldots,v_e).
\]
It is exactly \eqref{v4-arith:curve:eq:hall-split} at $s=-z$,
evaluated at $(x,y)=(u,v)$. Every finite auxiliary
derivative proves the asserted equality.

For each Laurent coefficient only a finite derivative
order of the polynomial test is needed, so its coefficient
belongs to the restricted dual. The preceding pole bound
gives a common lower bound on the fixed input.
Transposing $-\partial_x$ sends the $m$th derivative
functional to minus the $(m+1)$st one.
Finally $P$ is even, and
$r(-s)^{-1}=r(s)$ gives the stated exchange variance.
\end{proof}

With the opposite Hall grading,
$\deg(J_{m_1}\cdots J_{m_d})=2d^2-2\sum_i m_i$.
Translation has degree $-2$, and $\deg z=2$ makes the
transpose homogeneous. Ordinary multiplication in the
ambient Hopf algebra has a different charge-dependent
grading. Its cochain degree remains zero.

Charge completion gives $\prod_dW_{N,d}^+$ with outputs in
$\prod_dW_{N,d}^+((z))$. Each fixed total charge has finitely
many splittings, but there is no uniform pole bound over
all charges. The coefficient map
\eqref{v4-arith:curve:eq:hall-split} has not been extended
to every formal root series: translation by $s$ does not
preserve the root-adic topology.

The level map on the ambient Hopf carrier is different:
$(\theta_2\otimes1)\Delta$ sends
$G$ to $\theta_2(G)\otimes G$. Both output factors have
creation charge one. A Hall split has output charges
adding to its input charge. Moreover
$\vartheta(G)=(e^{\ell T_H}g)UV$ need not lie in $W_N^+$.

Thus the transpose theorem identifies Hall charge
cooperation with a field product on finite jets. It
does not identify it with a level-adding coproduct.
The corresponding field tensor, if that further comparison
is sought, must retain its spectral twist and this change
of charge.

\section{Moving centers and Cartan operations}
\label{v4-arith:curve:sec:moving-center}

Assume $N\geq4$ and extend the coefficients to
\[
 R=\mathbb C(a,b)[a_{\mathrm{fr}},\beta,\beta^{-1},
                    y_1,y_2,\kappa,\epsilon]/(\epsilon^N),
 \qquad \lambda=a_{\mathrm{fr}}+\epsilon b .
\]
The affine parameter $\kappa$ and labels $y_1,y_2$ are
retained independent coefficients. They are not the root
coordinates of the Hall correspondence.

On polynomial Hall coefficients define the connection
\begin{equation}\label{v4-arith:curve:eq:moving-connection}
 \nabla_d=\beta\left(\partial_{a_{\mathrm{fr}}}
                         +\sum_{i=1}^d\partial_{x_i}\right).
\end{equation}
On a tensor product over $R$ this is the tensor connection
over the single coefficient derivation
$\beta\partial_{a_{\mathrm{fr}}}$.
Every kernel depends only on root differences, so the
simultaneous root derivative kills it. Translation of the
first block by $s$ also commutes with this derivative.
Thus both the Hall product and
\eqref{v4-arith:curve:eq:hall-split} are horizontal.

The coordinates $\xi_i=x_i-\lambda$ are horizontal.
Completion along this moving section therefore has a
continuous dual with horizontal basis
\[
 v_{\mathbf m,\lambda}(F)
 =\left.\partial_{x_1}^{m_1}\cdots\partial_{x_d}^{m_d}
                         F\right|_{x_i=\lambda}.
\]
Indeed the chain rule gives
\[
 \beta\partial_{a_{\mathrm{fr}}}\langle v_{\mathbf m,\lambda},F\rangle
       =\langle v_{\mathbf m,\lambda},\nabla_dF\rangle .
\]

Sending this basis to $\prod_iJ_{m_i}$ gives the moving-center
version of the transpose theorem. Its proof evaluates $F$
at $(\lambda+z+u,\lambda+v)$, with the same scalar kernels.
The first insertion has the horizontal character
$F\mapsto F(\lambda)$. On $F=x$ both sides of the
horizontality identity are $\beta$.

Evaluation at a nonnilpotent $\lambda$ is not a finite
linear combination of derivatives at zero. These are jets
at a moving section, with no implicit $e^{\lambda T}$ on
the fixed algebraic state space. The state derivative $T$
raises jet order. The coefficient connection fixes the
displayed basis.

\subsection{The Cartan extension of the Hall coproduct}

Let $\mathsf H=R[\mathsf p_0,\mathsf p_1,\ldots]$ have
primitive generators acting on $A_d$ by multiplication
by $\sum_i x_i^n$. These act as derivations of the Hall
product, because the power sums add on concatenated root
lists. Form the smash product
$\mathsf B=(\bigoplus_dA_d)\#\mathsf H$ with relations
$[\mathsf p_n,F]=(\sum_i x_i^n)F$.
For a spectral variable $p$ define
\begin{equation}\label{v4-arith:curve:eq:cartan-current}
 \Phi(p)=\exp\left(\sum_{n\geq0}
       \frac{(-1)^n}{n!}\mathsf p_n\partial_p^n\log r(p)\right).
\end{equation}

Every coefficient in $p^{-1}$ uses finitely many
$\mathsf p_n$. On a finite root list $L$, substitution
$\mathsf p_n=\sum_{l\in L}l^n$ gives
$\Phi_L(p)=\prod_{l\in L}r(p-l)$.
This is a root-addition current. A framed highest-weight
current is further data.

\begin{proposition}\label{v4-arith:curve:prop:cartan-coproduct}
In the localized spectral tensor target, the formulas
\[
 \begin{aligned}
 \Delta^{\mathrm{ext}}_{d,e}(s)F
   &=\Delta^0_{d,e}(s)F\prod_{j=1}^e\Phi_L(y_j+s),\\
 \Delta^{\mathrm{ext}}(s)\mathsf p_n
   &=\sum_{i=0}^n\binom ni(-s)^{n-i}\mathsf p_i\otimes1
                                      +1\otimes\mathsf p_n
 \end{aligned}
\]
give a unital algebra coproduct on $\mathsf B$.
The coefficients of each $\Phi_L$ lie in the left output
factor, to the right of its Hall coefficient in the
displayed product. The coproduct is rationally coassociative
and horizontal when
$\nabla\mathsf p_n=\beta n\mathsf p_{n-1}$ and
$\nabla\mathsf p_0=0$.
\end{proposition}
\begin{proof}
For multiplicativity, split the first input into root
blocks $A,C$ and the second into $B,D$, with $A,B$ in
the left output. Cancellation of cross Euler factors leaves
\[
 \prod_{b\in B,c\in C}
      \frac{E_0(s+c-b)}{E_0(-s+b-c)}
 =\prod_{b\in B,c\in C}r(s+c-b).
\]

When the separate extended coproducts are multiplied,
the smash relation moves the first Cartan factor past
$B$ and replaces its list $L$ by $L\sqcup B$.
The factor changes by precisely this product. All other
factors coincide. The even factor $P$ cancels from the
ratio.

The binomial formula preserves the Cartan relations
with shifted roots.

Evaluation on all finite root lists is jointly faithful.
At length at least the largest positive power-sum index,
those power sums are algebraically independent.
Varying the length tests the polynomial in $\mathsf p_0$.
This remains valid coefficientwise modulo $\epsilon^N$.
The root calculation therefore proves the algebra identity.
For a triple split, both sides have Cartan factor
$\Phi_L(Y+s)\Phi_L(Z+s+t)\Phi_M(Z+t)$ and the same
Euler factors, proving coassociativity.

Finally differentiating \eqref{v4-arith:curve:eq:cartan-current}
by $\nabla\mathsf p_n=\beta n\mathsf p_{n-1}$ gives
$-\beta\partial_p\Phi$. It cancels the derivative
$\beta\partial_{y_j}\Phi_L(y_j+s)$.
The raw Euler part is horizontal, and the binomial formula
commutes with the same connection. This proves horizontality
on all generators.
\end{proof}

The Cartan operation also has a direct state realization.
For $C=G\vartheta(G)^{-1}$ set
\[
 \mathcal U(p)x=x_{(0)}\rho_p(C,x_{(1)}).
\]
It is an invertible algebra operator with
\[
 \mathcal U(p)G(v)=r(p-v)G(v),\qquad \mathcal U(p)1=1.
\]
Multiplicativity follows from the group-like $C$ and
the multiplicative coaction. $C^{-1}$ gives its inverse.
It preserves $W_N^+$ coefficientwise.

Pairing its products
and derivatives with polynomial tests shows that multiplication
by $\prod_i r(p-x_i)$ has finite-jet transpose
$\mathcal U(p-\lambda)$ at the moving center. In particular
$(\beta\partial_{a_{\mathrm{fr}}}+\beta\partial_p)
 \mathcal U(p-\lambda)=0$.

In the Laurent domain with $p$ large compared with $z$,
the creation field satisfies
\begin{equation}\label{v4-arith:curve:eq:cartan-covariance}
 \mathcal U(p)Y_\vartheta(G,z)\mathcal U(p)^{-1}
       =r(p-z)Y_\vartheta(G,z).
\end{equation}
Apply both sides to a product of $G(v_j)$.
The field adds the regular factor $G(z)$ and its
contractions. The Cartan operator multiplies this new
factor by $r(p-z)$ and cancels its factors on the input.
Differentiation proves the relation on all finite jets.
At $p-z=\beta u$, the cubic coefficient is
$2abh\beta^{-3}u^{-3}$.

\section{Framing markings and the supported current}
\label{v4-arith:curve:sec:framing}

The next calculation retains a moment complex with both
odd relation directions. Before localizing in $a,b$, use
\[
 R^0=\mathbb C[a,b,a_{\mathrm{fr}},\beta,\beta^{-1},
                   y_1,y_2,\kappa,\epsilon]/(\epsilon^N),
 \qquad N\geq4,\quad \lambda=a_{\mathrm{fr}}+\epsilon b,
 \quad c=ab(a+b).
\]
Its generic base change is the ring $R$ above.

The two odd directions occur already in a rank-two
framed moment chart. Take
\[
 X=\begin{pmatrix}x_1&B_0\\0&x_2\end{pmatrix},\quad
 Y=\begin{pmatrix}\upsilon_1&D_0\\0&\upsilon_2\end{pmatrix},
 \quad i=\binom{i_1}{i_2},\quad j=(1,0).
\]

The four entries of $[X,Y]+ij$ are
$i_1,(x_1-x_2)D_0-(\upsilon_1-\upsilon_2)B_0,i_2,0$.
Eliminating the two regular Koszul pairs for $i_1,i_2$
leaves odd generators $\eta,\chi$ with
\[
 d\eta=(x_1-x_2)D_0-(\upsilon_1-\upsilon_2)B_0,\qquad d\chi=0.
\]

The elimination is explicit: for $d\zeta=i$, the
contraction sends $i^m$ to $\zeta i^{m-1}$ for $m\geq1$.
Tensor the two contractions with their Koszul signs.

On the open $B_0\ne0$, the change of basis
$\operatorname{diag}(1,\gamma)$ preserves $j$ and the flag.
It scales $B_0,D_0$ by $\gamma^{-1}$, so uniquely fixes $B_0=1$.
Derived restriction to
$\upsilon_1=\upsilon_2=0$ then leaves $d\eta=uD_0$,
$u=x_1-x_2$, and still $d\chi=0$.
The zero-equation generator survives the restriction.

This specifies a moment coefficient chart. Orientations
for a Hall pushforward have not been supplied by the
matrix calculation.

In position coordinates $z,w,t$, put $u=z-w$, $v=t-w$,
and define
\begin{equation}\label{v4-arith:curve:eq:moment-module}
 \mathcal T=\bigoplus_{p,r\geq1}R^0[w]u^{-p}v^{-r},
 \qquad
 \mathcal M=\mathcal T[D_0]\otimes\Lambda(\eta,\chi),
 \quad d\eta=uD_0,\quad d\chi=0 .
\end{equation}

Here $\eta,\chi$ have degree $-1$, $D_0$ degree zero.
Principal-part multiplication discards nonnegative powers:
$u\,u^{-1}=v\,v^{-1}=0$ in their respective factors.
There is only a direct sum of normal orders.

For $\mathcal T_1=\ker(u:\mathcal T\to\mathcal T)$,
\begin{equation}\label{v4-arith:curve:eq:moment-cohomology}
 \begin{aligned}
 H^0(\mathcal M)&=\mathcal T,\\
 H^{-1}(\mathcal M)&=\eta\mathcal T_1[D_0]\oplus\chi\mathcal T,\\
 H^{-2}(\mathcal M)&=\eta\chi\mathcal T_1[D_0].
 \end{aligned}
\end{equation}

Indeed $u$ is surjective on $\mathcal T$, so the image
of $uD_0$ is $D_0\mathcal T[D_0]$ and its kernel is
$\mathcal T_1[D_0]$. Tensor with the closed $\chi$ factor.
The degree-zero and $\chi$ summands use $D_0=0$.

The required cubic contact class is
\begin{equation}\label{v4-arith:curve:eq:framing-contact}
 \gamma_N=2c\epsilon^3\beta^{-3}u^{-3}v^{-1}
                    \in H^0(\mathcal M).
\end{equation}
It is nonzero: $c$ is a non-zero-divisor in the polynomial
coefficient ring and $\epsilon^3\ne0$ at $N\geq4$.

\subsection{The three normal coefficients of a marking}

A scalar marking is an action of $R^0[\ell_{\mathrm{fr}}]$
on $\mathcal M$ in which $\ell_{\mathrm{fr}}$ multiplies by
$s(u,w)=\sum_{j\geq0}a_j(w)u^j$, with
$a_j\in R^0[w]$. Every principal part has finite $u$-pole
order, so this action is defined without completing
$\mathcal T$. It commutes with $d$.
It specifies a coefficient action, not the Hopf level
constructed in the preceding section.

\begin{proposition}\label{v4-arith:curve:prop:three-coefficients}
The relation $(s-\lambda)\gamma_N=0$ holds exactly when
\[
 a_0-\lambda,\quad a_1,\quad a_2
                    \ \in\epsilon^{N-3}R^0[w].
\]
At all compatible precisions, the coefficientwise criterion is
\[
 s=\lambda+u^3h(u,w).
\]
\end{proposition}
\begin{proof}
Multiplication by $u^j$ kills $u^{-3}$ for $j\geq3$, and
\[
 (s-\lambda)\gamma_N
 =2c\epsilon^3\beta^{-3}
       \bigl((a_0-\lambda)u^{-3}+a_1u^{-2}+a_2u^{-1}\bigr)v^{-1}.
\]
The three principal parts are independent. The annihilator
of $c\epsilon^3$ is $(\epsilon^{N-3})$, proving the criterion.
At all compatible precisions, each fixed $\epsilon^m$
coefficient of these three series vanishes by taking
$N>m+3$. The converse is immediate.
\end{proof}

In particular an unramified normal marking, for which
$a_1$ is a unit, cannot preserve the contact. Even
$s=\lambda+\beta u$, with the correct diagonal character,
leaves the defect $2c\epsilon^3\beta^{-2}u^{-2}v^{-1}$.
Modulo $\epsilon^4$ at $a=b=\beta=1$ its coefficient is $4$.
The ramified marking $s=\lambda+\beta u^3$ does preserve
the contact and gives the nonzero cochain map
$1\mapsto\gamma_N$ from $R^0[0]$ with character $\lambda$.
It has different translation behavior.

Let $\nabla=\partial_w$ at fixed $u,v$, fixing the
coefficient ring and the moment generators.
It commutes with $d$ and differentiates simultaneous
translation of $z,w,t$.
Call a marking translation-covariant with slope $\beta$
when $[\nabla,m_s]=\beta I$.

\begin{theorem}\label{v4-arith:curve:thm:marking-obstruction}
Such a marking has $s=\beta w+b(u)$ with
$b(u)\in R^0[[u]]$.
For every $\lambda'\in R^0$, multiplication by
$s-\lambda'$ is injective on every $H^i(\mathcal M)$.
Consequently a cochain map from a complex whose
cohomology has character $\lambda$ induces zero on
cohomology if it intertwines the marking, even up to
homotopy.
\end{theorem}
\begin{proof}
The commutator condition is $\partial_ws=\beta$.
Coefficientwise integration in characteristic zero gives
the displayed form.

On $\mathcal T$, take the largest
$u$-pole of a nonzero element. Its leading coefficient
is multiplied by $\beta w+b(0)-\lambda'$, a
non-zero-divisor because its leading coefficient in $w$
is the unit $\beta$. It remains nonzero.

On $\mathcal T_1[D_0]$, $u$ acts by zero and the same
polynomial proves injectivity. Formula
\eqref{v4-arith:curve:eq:moment-cohomology} treats the
closed odd factors as well.

For an intertwining map $\Phi$, compatibility on
cohomology gives
$(s-\lambda)H(\Phi)(x)=H(\Phi)((\ell_{\mathrm{fr}}-\lambda)x)=0$.
Injectivity forces $H(\Phi)(x)=0$.
A compatibility homotopy gives this same equality.
\end{proof}

Conjugating the marking by a cochain isomorphism or
adding $dH+Hd$ leaves the injectivity conclusion unchanged
on cohomology. The ramified marking above has
$[\nabla,m_s]=0$, so it is outside this theorem's
covariance hypothesis.

\subsection{The moving framing fiber}

Set
\[
 f=\beta(w+u)-\epsilon y_1-\lambda,\qquad
 \mathcal Q=\mathcal M/f\mathcal M.
\]
Multiplication by $f$ is injective on the terms and
cohomology of $\mathcal M$: its highest $w$ coefficient
is multiplied by the unit $\beta$.
Thus the ordinary quotient represents the derived
framing fiber and
$H^*(\mathcal Q)=H^*(\mathcal M)/fH^*(\mathcal M)$.
It eliminates $w$ by
$w=(\lambda+\epsilon y_1)/\beta-u$.
In degree zero the remaining basis is
$u^{-p}v^{-r}$ over $R^0$, and $\bar\gamma_N\ne0$.

The operator $\nabla=\partial_w$ does not descend:
$\nabla f=\beta$ would force $\mathcal Q=0$ if the
ideal were stable. Nor can an $R^0$-linear operator on
$\mathcal Q$ have commutator $\beta I$ with its marking,
which is the scalar $\lambda$.

The connection
\begin{equation}\label{v4-arith:curve:eq:framing-connection}
 \nabla_{\mathrm{fr}}=\partial_w+\beta\partial_{a_{\mathrm{fr}}}
\end{equation}
does descend. It annihilates $f$, commutes with the moment
differential, and satisfies
$[\nabla_{\mathrm{fr}},m_\lambda]=\beta I$.
Its coefficient rule is
$\nabla_{\mathrm{fr}}(rx)=\beta(\partial_{a_{\mathrm{fr}}}r)x+
r\nabla_{\mathrm{fr}}x$. This changes both the recipient and
the kind of connection. It does not contradict the
obstruction on the unchanged complex.

\subsection{The Cauchy comparison and both odd directions}

Work first over $A=R^0[z,w]=R^0[u,w]$.
For the current coordinate $t$ define
\[
 \begin{aligned}
 P_z&=A[t,(t-z)^{-1}]/A[t],&
 P_w&=A[t,(t-w)^{-1}]/A[t],\\
 P_{zw}&=A[t,((t-z)(t-w))^{-1}]/A[t],&
 \mathcal K_{zw}&=[P_z\oplus P_w\xrightarrow{(a,b)\mapsto a-b}P_{zw}].
 \end{aligned}
\]
The two terms have degrees $-1,0$ and use proper rational
representatives at $t=\infty$.
There is a relative differential-module quasi-isomorphism
$\Theta:\mathcal K_{zw}\to\mathcal T[0]$, zero in degree $-1$.
In degree zero, expand at $v=0$, take its $v$-principal
part, and then take its $u$-principal part.
Explicitly,
\begin{equation}\label{v4-arith:curve:eq:cauchy}
 \Theta\left(\frac1{v^a(v-u)^b}\right)
 =(-1)^b\sum_{j=0}^{a-1}
       \binom{b+j-1}{j}u^{-b-j}v^{j-a}.
\end{equation}

Here is a direct proof on the stated quotient. Monic
division by $(t-z)(t-w)$ makes $P_{zw}$ free over $A$.
$P_z,P_w$ have their usual inverse-power bases.
After inverting $u$, partial fractions separate the two
poles.

The integral summands have zero intersection.
Coefficientwise in $\epsilon$, a function belonging to
both has poles only at $z$ and only at $w$.
It therefore has no pole and, being proper, is zero.

Expansion gives
\eqref{v4-arith:curve:eq:cauchy} and kills both single-pole
summands. At fixed $a+b$, its matrix is triangular with
diagonal entries $(-1)^b$, so it surjects onto every
negative $u,v$ monomial.

If its value vanishes, the
$w$-principal part has no $u$-pole and lies in $P_w$.
Subtracting it leaves a member of $P_z$. This proves
the kernel assertion.

Principal parts commute with
derivatives, since differentiating a regular power
creates no pole. The physical derivatives are
$\partial_z=\partial_u$ and
$\partial_w^{\mathrm{phys}}=\partial_w-\partial_u-\partial_v$,
which proves relative differential linearity.

The derived diagonal restriction of $\mathcal K_{zw}$
has $P_w$ in degree $-1$ and zero in degree zero:
its differential becomes $(a,b)\mapsto a-b$ on two
copies of $P_w$. In contrast, the ordinary restriction
$\mathcal T/u\mathcal T$ is zero because $u$ is
surjective on $\mathcal T$.

Let $K_f=A[D_0]\otimes\Lambda(\eta,\chi)$ with
$d\eta=uD_0$, $d\chi=0$.
It is bounded with free $A$-terms, so
$1\otimes\Theta:K_f\otimes_A\mathcal K_{zw}\to\mathcal M$
is a quasi-isomorphism, fixing both odd generators.
The tensor differential
$d(\xi\otimes a)=d\xi\otimes a+(-1)^{|\xi|}\xi\otimes da$
and $A$-linearity prove its cochain equation.
It retains all the cohomology in
\eqref{v4-arith:curve:eq:moment-cohomology}.
It does not provide separate commuting position
connections on $K_f$: with $D_0$ fixed,
$[\partial_u,d]\eta=D_0$.

Now pass to the generic coefficient ring $R$ and the
current $G$ already constructed. Define
\[
 G_r=\frac{(\beta T)^{r-1}G}{(r-1)!},\qquad
 \mathcal J_G=\bigoplus_{r\geq1}RG_r,\qquad
 \mathcal Q_G=
 \left(\bigoplus_{p\geq1}Ru^{-p}\right)
       \otimes_R\mathcal J_G[D_0]\otimes\Lambda(\eta,\chi).
\]
Let $vG_1=0$, $vG_r=G_{r-1}$ for $r>1$, and let $w$
act by $(\lambda+\epsilon y_1)/\beta-u$.
Keep $d\eta=uD_0$, $d\chi=0$.
The coefficient connection is
$\beta\partial_{a_{\mathrm{fr}}}$, fixing this displayed basis.

\begin{theorem}\label{v4-arith:curve:thm:supported-current}
There is an isomorphism of coefficient complexes with
connection
\[
 \begin{aligned}
 \Psi:\mathcal Q\otimes_{R^0}R&\longrightarrow\mathcal Q_G,\\
 u^{-p}v^{-r}D_0^m\eta^e\chi^{e'}
       &\longmapsto u^{-p}G_rD_0^m\eta^e\chi^{e'}.
 \end{aligned}
\]
It retains both odd moment directions and satisfies
$\Psi(-\partial_v)=\beta T\Psi$.
The composition with $1\otimes\Theta$, followed by
derived base change to $f=0$, is a quasi-isomorphism
for the single common-translation connection
\eqref{v4-arith:curve:eq:framing-connection}.
It sends the cubic contact to
\[
 \Psi(\bar\gamma_N)=2abh\epsilon^3\beta^{-3}u^{-3}G_1.
\]
\end{theorem}
\begin{proof}
Elimination of $w$ gives a unique finite expression in
the displayed basis. The $G_r$ are independent, so $\Psi$ is an
isomorphism of modules. It preserves $u,D_0,\eta,\chi$
and the specified $v,w$ actions, hence the differential.
After eliminating $w$, $\nabla_{\mathrm{fr}}$ differentiates
only the coefficients of this basis, which proves
horizontality. The derivative identity follows from
$-\partial_v v^{-r}=rv^{-r-1}$ and
$\beta TG_r=rG_{r+1}$.

Simultaneous translation is
$\partial_z+\partial_w^{\mathrm{phys}}+\partial_t
=\partial_w$ at fixed $u,v$. It fixes $uD_0$.
The Cauchy map commutes with it and with the coefficient
derivative. $f$ is horizontal. Derived base change
preserves its quasi-isomorphism. Injectivity of $f$ on
terms and cohomology identifies its derived target with
the ordinary quotient $\mathcal Q$.

For the contact, expand the actual rational section
$(r(\beta(z-t))-1)(t-w)^{-1}$.
Its cubic coefficient is
$2abh\epsilon^3\beta^{-3}/((z-t)^3(t-w))$.
Formula~\eqref{v4-arith:curve:eq:cauchy}, with
$z-t=u-v$, gives the positive class
$2abh\epsilon^3\beta^{-3}u^{-3}v^{-1}$.
Applying $\Psi$ gives the stated value.
For $v^{-r}$, the same expansion retains the first
$r$ Taylor coefficients, precisely the corresponding
Leibniz derivatives of
\eqref{v4-arith:curve:eq:cartan-covariance}.
\end{proof}

Under the formal marking $x_i=\beta z_i-\epsilon y_i$,
the exact Cartan argument is
$\beta(z_1-z_2)-\epsilon(y_1-y_2)$.
Its correction first occurs at order $\epsilon^4$,
so the cubic normalization above is unchanged.
The common connection becomes
$\sum_i\partial_{z_i}+\beta\partial_{a_{\mathrm{fr}}}$,
and on the supported chart it is
\eqref{v4-arith:curve:eq:framing-connection}.
The separate derivative $\beta T$ differentiates the
relative output current.
These formulas do not identify plane eigenvalues in
the moment chart with tautological Chern roots of
a Hall correspondence.

The surviving classes include
$\eta u^{-1}G_1$ and $\eta\chi u^{-1}G_1$ in degrees
$-1,-2$. Hence the comparison cannot factor through
a degree-zero polynomial projection.
It also supplies only the common connection, with the
separate-position obstruction still given by
$[\partial_u,d]\eta=D_0$.

\section{The remaining geometric comparison}
\label{v4-arith:curve:sec:comparison}

The scalar charge-zero coaugmentation is $1$.
A local vacuum is a separate state marking. On the
test module $R\nu\otimes_R W_N^+$, take fields
$1\otimes Y_\vartheta$ and distinguished state
$\nu\otimes1$. Its first creation state is
$\nu\otimes G_1$, on which the first tautological
coefficient acts by $F(\lambda)$.
Then the identity field preserves $\nu\otimes1$.

This specifies the field test module and its framing
character. It does not identify it with a geometric
framed module.

\begin{problem}\label{v4-arith:curve:prob:critical}
Construct an oriented critical Hall/Hecke coefficient
family whose two-point map lands in
$K_f\otimes_A\mathcal K_{zw}$ and whose derived
framing fiber maps through
Theorem~\ref{v4-arith:curve:thm:supported-current}.
Require the Cartan and creation operations to give
\eqref{v4-arith:curve:eq:cartan-covariance}, with cubic
contact \eqref{v4-arith:curve:eq:framing-contact}.
Preserve the framing character $F(\lambda)$, both odd
moment directions, the common coefficient connection,
the scalar coaugmentation and the local vacuum marking.
Construct the comparison over $R$, and its descent to
$R^0$ before localization in $a,b$.

On the same coefficient family, construct the
nonconsecutive $13|2$ comparison and the required
small-diagonal homotopies. A comparison with a
level-adding coproduct must additionally specify its
different charge map and its twisted field tensor target.
\end{problem}

The zero-potential Euler calculation and the finite-jet
transpose are positive constructions on their named
carriers. The supported moving fiber also has a
constructed comparison, including its two odd directions
and one common connection. The oriented map from actual
critical Hall coefficients is still missing.
A physical identification must distinguish the nilpotent
coefficient parameter $\epsilon$ from a Moyal parameter
or a loop parameter.
It must also distinguish the exchange contact from a
gauge anomaly. Equal scalar
Euler factors alone give none of these further maps.
